% UTF-8 declarations (including ∈ for Lean code)


\documentclass[11pt,reqno]{amsart}
%%%%%%%%%%%%%%%%%%%%%%%%%%%%%%%%%%%%%%%%%%%%%%%%%%%%%%%%%%%%%%%%%%%%%%%%%%%%%%%%%%%%%%%%%%%%%%%%%%%%%%%%%%%%%%%%%%%%%%%%%%%%%%%%%%%%%%%%%%%%%%%%%%%%%%%%%%%%%%%%%%%%%%%%%%%%%%%%%%%%%%%%%%%%%%%%%%%%%%%%%%%%%%%%%%%%%%%%%%%%%%%%%%%%%%%%%%%%%%%%%%%%%%%%%%%%
\usepackage{amsfonts}
\usepackage[utf8]{inputenc}
\usepackage{amsmath,amssymb,amsthm}
\usepackage{changepage}
\usepackage{ragged2e}
\usepackage{float}
\usepackage{array}
\usepackage{graphicx}
\usepackage{setspace}
\usepackage{xcolor}
\usepackage{enumitem}
\usepackage[left=1in,right=1in,bottom=1in,top=1in,headheight=15pt,headsep=40pt]{geometry}
\usepackage{tikz}
\usepackage{fancyvrb}
\usepackage{fancyhdr}
\usepackage{titlesec,titletoc}
\usepackage[colorlinks=true,linkcolor=blue,citecolor=green,urlcolor=magenta]{hyperref}
\usepackage{aliascnt}
\usepackage[nameinlink]{cleveref}
\usepackage{listings}
\usepackage{booktabs}

\setcounter{MaxMatrixCols}{10}
%TCIDATA{OutputFilter=LATEX.DLL}
%TCIDATA{Version=5.50.0.2960}
%TCIDATA{<META NAME="SaveForMode" CONTENT="1">}
%TCIDATA{BibliographyScheme=Manual}
%TCIDATA{LastRevised=Saturday, November 14, 2026 19:08:23}
%TCIDATA{<META NAME="GraphicsSave" CONTENT="32">}

\restylefloat{figure}
\hypersetup{pageanchor=false}
\makeatletter
\@ifundefined{theorem}{}{\let\theorem\relax}
\@ifundefined{definition}{}{\let\definition\relax}
\@ifundefined{lemma}{}{\let\lemma\relax}
\@ifundefined{corollary}{}{\let\corollary\relax}
\@ifundefined{proposition}{}{\let\proposition\relax}
\@ifundefined{remark}{}{\let\remark\relax}
\@ifundefined{statement}{}{\let\statement\relax}
\makeatother
\DeclareUnicodeCharacter{2208}{\ensuremath{\in}}
\providecommand{\R}{\mathbb{R}}
\providecommand{\N}{\mathbb{N}}
\providecommand{\Q}{\mathbb{Q}}
\providecommand{\Z}{\mathbb{Z}}
\providecommand{\period}{\,\allowbreak}
\newcommand{\pagenumside}{C}
\newlength{\pageNumberOffset}\setlength{\pageNumberOffset}{0pt}
\newcommand{\PagenumSideC}{C}
\newcommand{\applypagehead}[1]{\expandafter\fancyhead\expandafter[\pagenumside]{#1}}
\newcommand{\applypagefoot}[1]{\expandafter\fancyfoot\expandafter[\pagenumside]{#1}}
\pagestyle{fancy}\fancyhf{}\applypagehead{\thepage}
\renewcommand{\headrulewidth}{0pt}
\renewcommand{\footrulewidth}{0pt}
\fancypagestyle{plain}{\fancyhf{}\applypagehead{\thepage}\renewcommand{\headrulewidth}{0pt}}
\setlength{\arraycolsep}{1.5pt}
\setlength{\tabcolsep}{3pt}
\renewcommand{\arraystretch}{1.05}
\setlength{\arrayrulewidth}{0.4pt}
\fvset{frame=lines,framerule=0.7pt,rulecolor=\color{blue!60!black},framesep=2pt,fontsize=\scriptsize,xleftmargin=4pt,xrightmargin=4pt}
\DefineVerbatimEnvironment{verbatimbox}{Verbatim}{frame=single,framerule=0.7pt,rulecolor=\color{blue!60!black},framesep=2pt,fontsize=\scriptsize,xleftmargin=4pt,xrightmargin=4pt}
\subjclass{\textbf{2020 Mathematics Subject Classification.}\,Primary 11A41, 11N05; Secondary 03F60, 03B35, 11-03.}
\keywords{prime-representing function, Mills' theorem, prime gaps, Ingham's estimate, Bertrand's postulate, constructive mathematics, Lean, verification}
\renewcommand{\abstractname}{\fontsize{14}{17}\selectfont Abstract}
\newcommand{\abstractfont}{\normalfont\fontsize{12}{14}\selectfont}
\newcommand{\abstractheadingsep}{\vspace{0.5em}}
\makeatletter
\renewenvironment{abstract}
 {\ifx\maketitle\relax\ClassWarning{\@classname}{Abstract should precede \protect\maketitle}\fi
  \global\setbox\abstractbox=\vtop \bgroup
  \abstractfont{\centering\abstractname\par}\abstractheadingsep
  \list{}{\labelwidth\z@\leftmargin2pc\rightmargin\leftmargin
         \listparindent\normalparindent\itemindent\z@\parsep\z@}
  \item\relax}
 {\endlist\egroup\ifx\@setabstract\relax\@setabstracta\fi}
\makeatother
\setcounter{secnumdepth}{3}
\setcounter{tocdepth}{1}
\newcommand{\SecAbove}{1.0ex plus 0.3ex minus 0.1ex}
\newcommand{\SecBelow}{0pt}
\newcommand{\SubsecAbove}{0.8ex plus 0.2ex minus 0.1ex}
\newcommand{\SubsecBelow}{0pt}
\newcommand{\SubsubsecAbove}{0.6ex plus 0.2ex minus 0.1ex}
\newcommand{\SubsubsecBelow}{0pt}
\makeatletter
\titleformat{\section}{\normalfont\fontsize{16}{18}\selectfont\bfseries}{\thesection.}{0.5em}{}
\titleformat{\subsection}{\normalfont\fontsize{14}{16}\selectfont\bfseries}{\thesubsection.}{0.5em}{}
\titleformat{\subsubsection}{\normalfont\fontsize{12}{14}\selectfont\bfseries}{\thesubsubsection.}{0.5em}{}
\titlespacing*{\section}{\z@}{\SecAbove}{\SecBelow}
\titlespacing*{\subsection}{\z@}{\SubsecAbove}{\SubsecBelow}
\titlespacing*{\subsubsection}{\z@}{\SubsubsecAbove}{\SubsubsecBelow}
\makeatother
\newcounter{axctr}\newcounter{posctr}\newcounter{defctr}
\newcounter{lemctr}\newcounter{thmctr}\newcounter{propctr}\newcounter{corctr}
\renewcommand{\theaxctr}{\arabic{axctr}}
\renewcommand{\theposctr}{\arabic{posctr}}
\renewcommand{\thedefctr}{\arabic{defctr}}
\renewcommand{\thelemctr}{\arabic{lemctr}}
\renewcommand{\thethmctr}{\arabic{thmctr}}
\renewcommand{\thepropctr}{\arabic{propctr}}
\renewcommand{\thecorctr}{\arabic{corctr}}
\renewcommand{\theHaxctr}{\thesection.\arabic{axctr}}
\renewcommand{\theHposctr}{\thesection.\arabic{posctr}}
\renewcommand{\theHdefctr}{\thesection.\arabic{defctr}}
\renewcommand{\theHlemctr}{\thesection.\arabic{lemctr}}
\renewcommand{\theHthmctr}{\thesection.\arabic{thmctr}}
\renewcommand{\theHpropctr}{\thesection.\arabic{propctr}}
\renewcommand{\theHcorctr}{\thesection.\arabic{corctr}}
\newtheoremstyle{redplain}{2pt plus 1pt minus 0.5pt}{1pt plus 0.5pt minus 0.3pt}{\itshape}{}{\bfseries\color{red}}{.}{\newline}{\thmname{#1}\thmnumber{ #2}\thmnote{ {\bfseries\color{red}(#3)}}}
\theoremstyle{redplain}
\newtheorem{theorem}{Theorem}[section]
\newtheorem{lemma}{Lemma}[section]
\newtheorem{corollary}{Corollary}[section]
\newtheorem{proposition}{Proposition}[section]
\newtheorem{remark}{Remark}[section]
\newtheorem{definition}{Definition}[section]
\newtheorem{axiom}{Axiom}[section]
\newtheorem{statement}{Statement}[section]
\definecolor{proofhead}{RGB}{11,83,148}
\definecolor{proofqed}{RGB}{204,0,0}
\renewcommand{\qedsymbol}{\textcolor{proofqed}{\ensuremath{\square}}}
\renewenvironment{proof}[1][Proof]{\par\noindent\textbf{\textcolor{proofhead}{#1.}}\par\nobreak}{\par\nobreak\hfill\qedsymbol}
\renewcommand{\theHtheorem}{\thesection.\arabic{theorem}}
\renewcommand{\theHstatement}{\thesection.\arabic{statement}}
\newcolumntype{p}[1]{>{\centering\arraybackslash\scriptsize}m{#1}|}
\crefname{theorem}{Theorem}{Theorems}
\Crefname{theorem}{Theorem}{Theorems}
\crefname{lemma}{Lemma}{Lemmas}
\Crefname{lemma}{Lemma}{Lemmas}
\crefname{axiom}{Axiom}{Axioms}
\Crefname{axiom}{Axiom}{Axioms}
\crefname{definition}{Definition}{Definitions}
\Crefname{definition}{Definition}{Definitions}
\crefname{remark}{Remark}{Remarks}
\Crefname{remark}{Remark}{Remarks}
\crefname{section}{Section}{Sections}
\Crefname{section}{Section}{Sections}
\crefname{table}{Table}{Tables}
\Crefname{table}{Table}{Tables}
\crefname{equation}{eq.}{eqs.}
\Crefname{equation}{Eq.}{Eqs.}
\newcounter{inq}
\renewcommand{\theinq}{\theequation}
\makeatletter\renewcommand{\p@inq}{\p@equation}\makeatother
\crefname{inq}{ineq.}{ineqs.}
\Crefname{inq}{Ineq.}{Ineqs.}
\makeatletter
\newenvironment{ineq}{\begin{equation}\refstepcounter{inq}}{\end{equation}}
\makeatother
\makeatletter
\let\orig@maketitle\maketitle
\renewcommand{\maketitle}{\begingroup\let\MakeUppercase\@firstofone\let\uppercase\@firstofone\orig@maketitle\endgroup}
\newlength{\titleToAuthorSkip}\newlength{\authorToAbstractSkip}
\setlength{\titleToAuthorSkip}{4pt plus 1pt minus 0.5pt}
\setlength{\authorToAbstractSkip}{6pt plus 1pt minus 0.5pt}
\renewcommand{\@maketitle}{\normalfont\normalsize\let\@makefnmark\relax\let\@thefnmark\relax
 \ifx\@empty\@subjclass\else\@footnotetext{\@setsubjclass}\fi
 \ifx\@empty\@keywords\else\@footnotetext{\@setkeywords}\fi
 \@settitle
 \ifx\@empty\@authors\else\vspace{\titleToAuthorSkip}\@setauthors\fi
 \ifx\@empty\@abstract\else\vspace{\authorToAbstractSkip}\fi
 \@setabstract}
\makeatother
\renewcommand{\encodingdefault}{OT1}
\renewcommand{\rmdefault}{cmr}
\renewcommand{\sfdefault}{cmss}
\renewcommand{\ttdefault}{cmtt}
\fontencoding{OT1}\fontfamily{cmr}\selectfont
\renewcommand{\contentsname}{\texorpdfstring{{\color{blue!60!black}\sffamily\Large\selectfont\bfseries Contents}}{Contents}}
\renewcommand{\seriesdefault}{\bfdefault}
\boldmath
\contentsmargin{1em}
\titlecontents{part}[0pt]{\addvspace{14pt}\large\bfseries}{\makebox[2.0em][l]{\thecontentslabel.}}{}{\titlerule*[0.6pc]{.}\contentspage}[\addvspace{4pt}]
\titlecontents{section}[2.0em]{\addvspace{10pt}\small\bfseries}{\makebox[2.0em][l]{\thecontentslabel.}}{\hspace*{-2.0em}}{\titlerule*[0.6pc]{.}\contentspage}[\addvspace{2pt}]
\titlecontents{subsection}[4.5em]{\addvspace{4pt}\footnotesize}{\makebox[2.5em][l]{\thecontentslabel.}}{\hspace*{-2.5em}}{\titlerule*[0.6pc]{.}\contentspage}[\addvspace{1pt}]
\setlist[itemize]{topsep=0pt,partopsep=0pt,itemsep=1pt plus 0.5pt minus 0.3pt,parsep=0pt,leftmargin=1.5em,labelsep=0.4em,label=\textbullet}
\setlist[enumerate]{topsep=0pt,partopsep=0pt,itemsep=1pt plus 0.5pt minus 0.3pt,parsep=0pt,leftmargin=1.8em,labelsep=0.4em}
\setlength{\parindent}{0pt}
\setlength{\parskip}{0.25em plus 0.05em minus 0.05em}
\makeatletter
\renewenvironment{quotation}{\list{}{\listparindent\z@\itemindent\z@\rightmargin\leftmargin\parsep\z@}\item\relax}{\endlist}
\makeatother
\numberwithin{equation}{section}
\lstdefinestyle{python}{
  language=Python,
  basicstyle=\ttfamily\scriptsize,
  keywordstyle=\color{blue},
  commentstyle=\color{gray},
  stringstyle=\color{orange},
  numbers=left,
  numberstyle=\tiny,
  breaklines=true,
  frame=single,
  showstringspaces=false,
  xleftmargin=4pt,
  xrightmargin=4pt
}
\IfFileExists{tcilatex.tex}{\input{tcilatex}}{}
\addtocontents{toc}{\setcounter{tocdepth}{1}}

\begin{document}
\title{\bfseries{\huge \emph{Mills' Theorem.}}\\
[12pt]{\Large \emph{From Ingham's Prime-Gap Estimate to Lean's Formulation}} 
\\
[6pt]{\Large \emph{Ten Routes to a Prime-Representing Function}}}
\author{ 
\bfseries{\large \hspace{1.2in} Brahimi Mahdi Tahar\newline\text{ \ 
			} \newline\text{ \  } \hspace{0.5in}\text{ \  } University of Mohamed
			Boudiaf of M'sila, Algeria\newline\text{ \  }
			\newline\href{mailto:mahditahar.brahimi@univ-msila.dz}{mahditahar.brahimi@univ-msila.dz}}
}
\maketitle

\begin{abstract}
\thispagestyle{empty} \noindent \setstretch{1.2}We collect ten routes to the
prime-representing theorem of W.~H.~Mills (1947): 
\begin{equation}
\exists A>1,\ \forall n\in \N_{\geq 1}:\quad \lfloor A^{3^{n}}\rfloor \text{
is prime.}
\end{equation}

\begin{itemize}
\item \noindent Three are prime-gap constructions: Mills' original proof
(via Ingham's \cite{Ingham1937} bound 
\begin{equation}
p_{n+1}-p_{n}<k\,p_{n}^{5/8}),
\end{equation}
Kuipers' generalization to $c^{n}$ for integer $c\geq 3$, and Ansari's
generalization to $c>8/3$.

\item Three are elementary or abstract: Wright's $2$-tower construction via
Bertrand's postulate \cite{HardyWright1960}, Ore's unified
selection-function framework, and a direct construction from a general
polynomial prime-gap estimate.

\item Two are variants and metrical: Niven's $c^{\theta^n}$ variant and
Wright's metrical theorem on the set of suitable $\theta$.

\item \noindent Two concern the Lean formulation: a pointer to the Lean
formalization in \texttt{formal-conjectures} \cite{FormalConjectures}, and a
detailed mathematical adaptation of the Lean proof using the \texttt{IsMills}
predicate and the sequence-limit construction to extract $\theta$ from the
recurrence.
\end{itemize}

\noindent The prime-gap input (Ingham \cite{Ingham1937}, 1937) is used
throughout: 
\begin{equation}
\forall n\in \N_{\geq 1}:\quad p_{n+1}-p_{n}<k\,p_{n}^{5/8}.
\end{equation}

\noindent The \emph{theorem} is $\mathsf{RCA}_{0}$-provable from the
recurrence, under Bishop--Cauchy coding of the real $\theta$ as a modulated
Cauchy sequence of rationals.

\noindent The \emph{routes} that invoke Ingham's estimate (Routes 1, 2, 3,
and 7) are provable in $\mathsf{ZF}$ but are \emph{not known to be} $\mathsf{%
RCA}_{0}$ -provable unconditionally, since Ingham's estimate itself is not
known to lie within $\mathsf{RCA}_{0}$ or $\mathsf{WKL}_{0}$. Wright's
Bertrand-based route (Route 4) sits at $\mathsf{RCA}_{0}$, since Bertrand's
postulate is elementary.

\noindent Route 6, the direct construction, inherits the strength of the
general polynomial prime-gap estimate that is fed into it; with Bertrand's
postulate, only $2$-towers are reachable, so the $3^n$ statement requires
Ingham's estimate.
\end{abstract}

\tableofcontents

\newpage 
\hypersetup{pageanchor=false} {\ \thispagestyle{empty} 
\pagestyle{empty}   \thispagestyle{empty}  \pagestyle{empty} } \newpage

\section{Introduction}

\setstretch{1.2} \pagestyle{plain} \pagenumbering{arabic} %
\hypersetup{pageanchor=true} \label{sec:introduction}

The prime-representing function of W.~H.~Mills~\cite{Mills1947} is one of
the more striking results of twentieth-century number theory: there exists a
real $\theta$ such that $[\theta^{3^n}]$ is prime for every $n=1,2,\dots$.
The formula yields prime values exclusively, and yet it is not
computationally useful: to construct $\theta$ one must already be able to
recognize arbitrarily large primes.

\medskip The purpose of this collection is to record ten routes to the same
conclusion, each independent and each verifying the primality of $%
[\theta^{3^n}]$ by different means. The geometric content of the recurrence
is shown in \Cref{fig:mills-recurrence}.

The Mills recurrence is

\begin{ineq}
\label{eq:mills-seq-intro} q_n^3<q_{n+1}<(q_n+1)^3-1\qquad(n=1,2,\dots),
\end{ineq}

and the limit $\theta $ is obtained as the common limit of 
\begin{equation}
u_{n}=q_{n}^{3^{-n}}\quad \text{and}\quad v_{n}=(q_{n}+1)^{3^{-n}}.
\end{equation}

\begin{figure}[]
\centering
\begin{tikzpicture}[xscale=0.55,yscale=0.9]
\draw[->] (0,0) -- (11,0) node[right] {$n$};
\draw[->] (0,0) -- (0,5) node[above] {$\log q_n$};
\draw[blue,thick] (0.5,0.5) -- (3,2.0) -- (6,3.6) -- (10,4.6);
\fill (0.5,0.5) circle (1.5pt) node[below right, inner sep=2pt] {$q_1$};
\fill (3,2.0)   circle (1.5pt) node[below right, inner sep=2pt] {$q_2$};
\fill (6,3.6)   circle (1.5pt) node[below right, inner sep=2pt] {$q_3$};
\fill (10,4.6)  circle (1.5pt) node[below left,  inner sep=2pt] {$q_4$};
\draw[dashed] (0.5,0) -- (0.5,0.5);
\draw[dashed] (3,0)   -- (3,2.0);
\draw[dashed] (6,0)   -- (6,3.6);
\draw[dashed] (10,0)  -- (10,4.6);
\node[blue,anchor=west] at (11,3.5) {$\log q_n$};
\end{tikzpicture}
\caption{The Mills recurrence: $\log q_n$ grows by a factor $3$ at each
step, so $\log\protect\theta=\lim 3^{-n}\log q_n$. (The printed $5^{-n}$ in
the historical transcription is a typo; see \Cref{rem:typo}.)}
\label{fig:mills-recurrence}
\end{figure}

\noindent The plan is as follows:

\begin{itemize}
\item \Cref{sec:prime-gap} collects the three prime-gap constructions
(Mills, Kuipers, Ansari).

\item \Cref{sec:elementary} collects the three elementary or abstract
constructions (Wright via Bertrand, Ore's framework, a direct route).

\item \Cref{sec:variants} collects three further routes (Niven's variant,
Wright's metrical theorem, and a pointer to Lean).

\item \Cref{sec:lean} reproduces the Lean formulation and proves the
associated statements.

\item \Cref{sec:python} presents the Python verification script.

\item \Cref{sec:tables} collects the verification tables for $N=20$.

\item \Cref{sec:summary} re-classifies the ten routes in second-order
arithmetic.
\end{itemize}

\begin{definition}[Mills sequence, primality predicate]
\label{def:mills-seq} Let $\{q_{n}\}_{n\geq 1}$ be a sequence of primes
satisfying 
\begin{equation}
\forall n\in 
%TCIMACRO{\U{2115} }%
%BeginExpansion
\mathbb{N}
%EndExpansion
_{\geq 1}\text{: \ \ \ }q_{n}^{3}<q_{n+1}<(q_{n}+1)^{3}-1
\end{equation}
for all $n\geq 1$. We write $\mathrm{Mills}(q)$ for the assertion that such
a sequence exists, and $\theta $ for its limit constant.
\end{definition}

\begin{definition}[Mills predicate, informal]
\label{def:ismills-intro} For a real $A$, the assertion $\mathrm{IsMills}(A)$
means that $\mathrm{Prime}(\lfloor A^{3^n}\rfloor)$ holds for every $n\ge1$,
where $\lfloor\cdot\rfloor$ denotes the integer-part function. Note that $%
\mathrm{IsMills}(A)$ does \emph{not} require $A^{3^n}$ itself to be an
integer; only its floor is required to be prime. The constraint used in
Mills' theorem is the joint condition $A>1$ together with $\mathrm{IsMills}%
(A)$.
\end{definition}

\begin{definition}[Cauchy codes, base theories]
\label{def:cauchy} A fast Cauchy code for a real is a sequence $q:\N%
\rightarrow \Q$ with $|$%
\begin{equation}
q_{m}-q_{n}|\leq 2^{-n}
\end{equation}%
for $m\geq n$. Two such codes represent the same real if their difference
tends to zero. Set 
\begin{equation}
\begin{aligned} \mathsf{RCA}_0&:=\text{basic second-order arithmetic}
+\Delta^0_1\text{-comprehension}+\Sigma^0_1\text{-induction},\\
\mathsf{WKL}_0&:=\mathsf{RCA}_0+\mathsf{WKL}, \end{aligned}
\end{equation}%
where $\mathsf{WKL}$ asserts every infinite binary tree has an infinite path.

%TCIMACRO{\TeXButton{noindent}{\noindent}}%
%BeginExpansion
\noindent%
%EndExpansion
Real quantifiers in coded statements range over valid Cauchy codes.
\end{definition}

%=========================================================================

\section{Prime-gap constructions}

\label{sec:prime-gap}

\subsection{Mills' original proof (First proof)}

\label{sec:mills}

\begin{theorem}[Mills \protect\cite{Mills1947}]
\label{thm:mills} There is a real number $\theta$ such that $[\theta^{3^n}]$
is prime for every $n=1,2,\dots$.
\end{theorem}

\begin{proof}
\noindent By Ingham's estimate \cite{Ingham1937}, 
\begin{equation}
\text{there is }k>0\text{ with }p_{n+1}-p_{n}<k\,p_{n}^{5/8}\text{ for all }%
n\geq 1.\text{ }
\end{equation}
\medskip

\noindent 
\begin{tabular}{|@{}p{0.06\linewidth}p{0.61\linewidth}p{0.29\linewidth}}
\hline
1 & $k>0$, \ $p_{n+1}-p_n<k\,p_n^{5/8}$ for all $n$ & Ingham's estimate \cite%
{Ingham1937}. \\ \hline
\medskip 2 & \medskip $N>k^8$, \ $p_n$ largest prime $<N^3$ & Arbitrary
parameters. \\ 
&  & 
\end{tabular}%
\noindent \text{ \ }\newline
\noindent 
\begin{tabular}{|@{}p{0.06\linewidth}p{0.61\linewidth}p{0.29\linewidth}}
\hline
3 & $p_n<N^3<p_{n+1}<N^3+N^2<(N+1)^3-1$ & Arithmetic; 1--2. \\ \hline
\medskip 4 & \medskip $\exists$ prime in $(N^3,(N+1)^3-1)$ for every $N>k^8$
& Existential introduction; 3. \\ 
&  & 
\end{tabular}%
\noindent \text{ \ }\newline
\noindent 
\begin{tabular}{|@{}p{0.06\linewidth}p{0.61\linewidth}p{0.29\linewidth}}
\hline
5 & $q_1>k^8$ prime; \ $q_{n+1}$ smallest prime in $(q_n^3,(q_n+1)^3-1)$ & 
Definite description (least element); 4. \\ \hline
\medskip 6 & \medskip $q_n^3<q_{n+1}<(q_n+1)^3-1$ for all $n\ge1$ & 
Recurrence; 5. \\ 
&  & 
\end{tabular}%
\noindent \text{ \ }\newline
\noindent 
\begin{tabular}{|@{}p{0.06\linewidth}p{0.61\linewidth}p{0.29\linewidth}}
\hline
7 & $u_n:=q_n^{3^{-n}}$, \ $v_n:=(q_n+1)^{3^{-n}}$ & Definition. \\ \hline
\medskip 8 & \medskip $u_n<u_{n+1}$ \ and \ $v_{n+1}<v_n$ for all $n\ge1$ & 
Cubing; 6--7. \\ 
&  & 
\end{tabular}%
\noindent \text{ \ }\newline
\noindent 
\begin{tabular}{|@{}p{0.06\linewidth}p{0.61\linewidth}p{0.29\linewidth}}
\hline
9 & $u_n<v_n$ for all $n\ge1$ & Algebra; 7. \\ \hline
\medskip 10 & \medskip $u_n<u_{n+1}<v_{n+1}<v_n$ for all $n\ge1$ & Chain;
8--9. \\ 
&  & 
\end{tabular}%
\noindent \text{ \ }\newline
\noindent 
\begin{tabular}{|@{}p{0.06\linewidth}p{0.61\linewidth}p{0.29\linewidth}}
\hline
11 & $\theta:=\lim u_n=\lim v_n$ exists & Cauchy with modulus; 10. \\ \hline
\medskip 12 & \medskip $q_n<\theta^{3^n}<q_n+1$ for all $n\ge1$ & $3^n$-th
power; 7, 11. \\ 
&  & 
\end{tabular}%
\noindent \text{ \ }\newline
\noindent 
\begin{tabular}{|@{}p{0.06\linewidth}p{0.61\linewidth}p{0.29\linewidth}}
\hline
13 & $[\theta^{3^n}]=q_n$ is prime for all $n\ge1$ & Floor; 12. \\ \hline
\medskip 14 & \medskip $\top$ & $\rightarrow$I; 1--13. \\ \hline
\end{tabular}
\end{proof}

\begin{remark}[Transcription typo]
\label{rem:typo} The historical transcription of Dudley \cite{Dudley1969}
writes 
\begin{equation}
u_{n}=q_{n}^{5^{-n}},v_{n}=(q_{n}+1)^{5^{-n}},
\end{equation}
and concludes%
\begin{equation}
q_{n}=[\theta ^{5^{n}}].
\end{equation}

%TCIMACRO{\TeXButton{noindent}{\noindent}}%
%BeginExpansion
\noindent%
%EndExpansion
That is a typo.

With exponent $5^{-n}$ the sequence $\{u_{n}\}$ fails monotonicity at $n=1$:

since 
\begin{equation}
q_{n+1}>q_{n}^{3}
\end{equation}%
and 
\begin{equation}
3\cdot 5^{-(n+1)}=3/5\cdot 5^{-n}<5^{-n},
\end{equation}%
we get 
\begin{equation}
u_{n+1}<u_{n}.
\end{equation}
Only the exponent $3^{-n}$ (matching the recurrence \Cref{eq:mills-seq-intro}%
) restores monotonicity. See \Cref{sec:tables} for a numerical confirmation.
\end{remark}

\subsection{Kuipers' generalization to $c^n$ (Second proof)}

\label{sec:kuipers}

\begin{theorem}[Kuipers \protect\cite{Kuipers1950}]
\label{thm:kuipers} If $c\ge3$ is an integer, then there is a real number $%
\theta$ such that $[\theta^{c^n}]$ is prime for all $n=1,2,\dots$.
\end{theorem}

\begin{proof}
\begin{tabular}{|@{}p{0.06\linewidth}p{0.61\linewidth}p{0.29\linewidth}}
\hline
1 & $c\ge3$ integer; \ $a=3c-4$, \ $b=3c-1$, \ $a/b\ge5/8$ & Hypothesis;
algebra. \\ \hline
\medskip 2 & \medskip $k_1>0$, \ $p_{n+1}-p_n<k_1\,p_n^{a/b}$ for all $n$ & 
Ingham \cite{Ingham1937}; 1. \\ 
&  & 
\end{tabular}%
\noindent \text{ \ }\newline
\noindent 
\begin{tabular}{|@{}p{0.06\linewidth}p{0.61\linewidth}p{0.29\linewidth}}
\hline
3 & $N>k_1^b$, \ $p_n$ largest prime $<N^c$ & Arbitrary parameters. \\ \hline
\medskip 4 & \medskip $p_n<N^c<p_{n+1}<N^c+N^{c-1}<(N+1)^c-1$ & Arithmetic; $%
ca+1=(c-1)b$; 1--3. \\ 
&  & 
\end{tabular}%
\noindent \text{ \ }\newline
\noindent 
\begin{tabular}{|@{}p{0.06\linewidth}p{0.61\linewidth}p{0.29\linewidth}}
\hline
5 & $\exists$ prime in $(N^c,(N+1)^c-1)$ for every $N>k_1^b$ & Existential;
4. \\ \hline
\medskip 6 & \medskip $q_n^c<q_{n+1}<(q_n+1)^c-1$; \ $u_n$, $v_n$, $\theta$
as in \Cref{thm:mills} & Recurrence; 5. \\ 
&  & 
\end{tabular}%
\noindent \text{ \ }\newline
\noindent 
\begin{tabular}{|@{}p{0.06\linewidth}p{0.61\linewidth}p{0.29\linewidth}}
\hline
7 & $q_n=[\theta^{c^n}]$ for all $n\ge1$ & Floor; 6. \\ \hline
\medskip 8 & \medskip $\top$ & $\rightarrow$I; 1--7. \\ \hline
\end{tabular}
\end{proof}

\subsection{Ansari's generalization to $c>8/3$ (Third proof)}

\label{sec:ansari}

\begin{theorem}[Ansari \protect\cite{Ansari1951}]
\label{thm:ansari} If $c>8/3$, there is a real number $\theta$ such that $%
[\theta^{c^n}]$ is prime for all $n=1,2,\dots$.
\end{theorem}

\begin{proof}
\begin{tabular}{|@{}p{0.06\linewidth}p{0.61\linewidth}p{0.29\linewidth}}
\hline
1 & $c>8/3$, \ $\delta=3c/8-1>0$ & Hypothesis; algebra. \\ \hline
\medskip 2 & \medskip $k>0$, \ $p_{n+1}-p_n<k\,p_n^{5/8}$ for all $n$ & 
Ingham \cite{Ingham1937}. \\ 
&  & 
\end{tabular}%
\noindent \text{ \ }\newline
\noindent 
\begin{tabular}{|@{}p{0.06\linewidth}p{0.61\linewidth}p{0.29\linewidth}}
\hline
3 & $N>k^{1/\delta}$, \ $p_n$ largest prime $<N^c$ & Arbitrary parameters.
\\ \hline
\medskip 4 & \medskip $p_n<N^c<p_{n+1}<N^c+N^{c-1}<(N+1)^c-\delta$ & 
Arithmetic; 1--3. \\ 
&  & 
\end{tabular}%
\noindent \text{ \ }\newline
\noindent 
\begin{tabular}{|@{}p{0.06\linewidth}p{0.61\linewidth}p{0.29\linewidth}}
\hline
5 & $\exists$ prime in $(N^c,(N+1)^c-1)$ for every $N>k^{1/\delta}$ & 
Existential; 4. \\ \hline
\medskip 6 & \medskip $q_n^c<q_{n+1}<(q_n+1)^c-1$; \ $\theta$ as in %
\Cref{thm:kuipers} & Recurrence; 5. \\ 
&  & 
\end{tabular}%
\noindent \text{ \ }\newline
\noindent 
\begin{tabular}{|@{}p{0.06\linewidth}p{0.61\linewidth}p{0.29\linewidth}}
\hline
7 & $q_n=[\theta^{c^n}]$ for all $n\ge1$ & Floor; 6. \\ \hline
\medskip 8 & \medskip $\top$ & $\rightarrow$I; 1--7. \\ \hline
\end{tabular}
\end{proof}

\begin{remark}
Ansari's exponent $8/3$ can be improved.

%TCIMACRO{\TeXButton{noindent}{\noindent}}%
%BeginExpansion
\noindent%
%EndExpansion
Titchmarsh showed that $5/8$ in Ingham's estimate could be replaced by $%
48/77 $, giving 
\begin{equation}
c>77/29.
\end{equation}%
%TCIMACRO{\TeXButton{noindent}{\noindent}}%
%BeginExpansion
\noindent%
%EndExpansion
Haneke's estimate $61/98+\epsilon $ gives 
\begin{equation}
c>98/37.
\end{equation}
\end{remark}

%=========================================================================

\section{Elementary and abstract constructions}

\label{sec:elementary}

\subsection{Wright's $2$-tower via Bertrand (Fourth proof)}

\label{sec:wright}

\begin{theorem}[Wright \protect\cite{Wright1951}]
\label{thm:wright} There is a real number $\theta $ such that $[2^{2^{\cdot
^{\cdot ^{2^{\theta }}}}}]$ (with $n$ iterations of the exponential) is
prime for every $n\geq 1$.%
\begin{equation*}
\end{equation*}
\end{theorem}

\begin{proof}
\begin{tabular}{|@{}p{0.06\linewidth}p{0.61\linewidth}p{0.29\linewidth}}
\hline
1 & $\forall N\ge2$, \ $\exists$ prime in $(N,2N)$ & Bertrand \cite%
{HardyWright1960}. \\ \hline
\medskip 2 & \medskip $p_n<N<p_{n+1}<2p_n<2N$ for $p_n$ largest prime $<N$ & 
Order; 1. \\ 
&  & 
\end{tabular}%
\noindent \text{ \ }\newline
\noindent 
\begin{tabular}{|@{}p{0.06\linewidth}p{0.61\linewidth}p{0.29\linewidth}}
\hline
3 & $q_1$ prime; \ $q_{n+1}$ smallest prime in $(2^{q_n},2^{q_n+1})$ & 
Definite description (least element); 2. \\ \hline
\medskip 4 & \medskip $q_n<\log q_{n+1}<\log(q_{n+1}+1)<q_n+1$ & Base-$2$
log; 3. \\ 
&  & 
\end{tabular}%
\noindent \text{ \ }\newline
\noindent 
\begin{tabular}{|@{}p{0.06\linewidth}p{0.61\linewidth}p{0.29\linewidth}}
\hline
5 & $u_n:=\log^{(n)}q_n$, \ $v_n:=\log^{(n)}(q_n+1)$ & Definition. \\ \hline
\medskip 6 & \medskip $u_n<u_{n+1}<v_{n+1}<v_n$ for all $n\ge1$ & Iterated
log; 4--5. \\ 
&  & 
\end{tabular}%
\noindent \text{ \ }\newline
\noindent 
\begin{tabular}{|@{}p{0.06\linewidth}p{0.61\linewidth}p{0.29\linewidth}}
\hline
7 & $\theta:=\lim u_n=\lim v_n$ exists & Cauchy; 6. \\ \hline
\medskip 8 & \medskip $q_n<\exp^{(n)}\theta<q_n+1$ for all $n\ge1$ & 
Exponentiate; 5, 7. \\ 
&  & 
\end{tabular}%
\noindent \text{ \ }\newline
\noindent 
\begin{tabular}{|@{}p{0.06\linewidth}p{0.61\linewidth}p{0.29\linewidth}}
\hline
9 & $[\exp^{(n)}\theta]=q_n$ is prime for all $n\ge1$ & Floor; 8. \\ \hline
\medskip 10 & \medskip $\top$ & $\rightarrow$I; 1--9. \\ \hline
\end{tabular}
\end{proof}

\begin{remark}
Only the $2$-tower is reachable from Bertrand. The $3^n$ version of %
\Cref{thm:mills} requires Ingham's estimate, which is not elementary.
\end{remark}

\subsection{Ore's unified framework (Fifth proof)}

\label{sec:ore}

\begin{theorem}[Ore \protect\cite{Ore1952}]
\label{thm:ore} Let $f$ be a selection function for a prime sequence $%
\{q_{n}\}$ continuous, eventually increasing, and such that 
\begin{equation*}
f(q_{n}-\delta _{n})<q_{n+1}-\delta _{n+1}\leq q_{n+1}+\epsilon
_{n+1}<f(q_{n}+\epsilon _{n})
\end{equation*}
for some $\delta _{n},\epsilon _{n}\geq 0$). Then 
\begin{equation*}
q_{n}=[f^{(n)}(\theta )]
\end{equation*}%
for a suitable limit $\theta $.
\end{theorem}

\begin{proof}
\begin{tabular}{|@{}p{0.06\linewidth}p{0.61\linewidth}p{0.29\linewidth}}
\hline
1 & $f$ continuous, eventually increasing; \ $f(q_n-\delta_n)<q_{n+1}-%
\delta_{n+1}\le q_{n+1}+\epsilon_{n+1}<f(q_n+\epsilon_n)$ & Selection
hypothesis. \\ \hline
\medskip 2 & \medskip $u_n:=f^{(-n)}(q_n-\delta_n)$, \ $v_n:=f^{(-n)}(q_n+%
\epsilon_n)$ & Definition. \\ 
&  & 
\end{tabular}%
\noindent \text{ \ }\newline
\noindent 
\begin{tabular}{|@{}p{0.06\linewidth}p{0.61\linewidth}p{0.29\linewidth}}
\hline
3 & $u_n<u_{n+1}<v_{n+1}<v_n$ & Apply $f^{(-(n+1))}$; 1--2. \\ \hline
\medskip 4 & \medskip $\theta:=\lim u_n=\lim v_n$ & Cauchy; 3. \\ 
&  & 
\end{tabular}%
\noindent \text{ \ }\newline
\noindent 
\begin{tabular}{|@{}p{0.06\linewidth}p{0.61\linewidth}p{0.29\linewidth}}
\hline
5 & $q_n-\delta_n\le f^{(n)}(\theta)\le q_n+\epsilon_n$ for all $n\ge1$ & 
Apply $f^{(n)}$; 1, 4. \\ \hline
\medskip 6 & \medskip $q_n=[f^{(n)}(\theta)]$ when $\delta_n=\epsilon_n=0$ & 
Floor; 5. \\ 
&  & 
\end{tabular}%
\noindent \text{ \ }\newline
\noindent 
\begin{tabular}{|@{}p{0.06\linewidth}p{0.61\linewidth}p{0.29\linewidth}}
\hline
7 & $\top$ & $\rightarrow$I; 1--6. \\ \hline
\end{tabular}
\end{proof}

\begin{remark}
Ore's framework unifies \Cref{thm:mills}, \Cref{thm:kuipers}, %
\Cref{thm:ansari}, and \Cref{thm:wright}:

\begin{enumerate}
\item Mills selection function is%
\begin{equation}
f(x)=x^{3}
\end{equation}

\item Kuiperss selection function is 
\begin{equation}
f(x)=x^{c}
\end{equation}

\item Ansari; $f(x)=2^{x}$ gives Wright.
\end{enumerate}
\end{remark}

\subsection{Direct construction from a polynomial gap (Sixth proof)}

\label{sec:direct}

\begin{theorem}
\label{thm:direct} Suppose there exist $k>0$ and $\alpha\in(0,1)$ with $%
p_{n+1}-p_n<k\,p_n^{\alpha}$ for all $n$. Then there is a real $\theta$ such
that $[\theta^{(1/\alpha)^n}]$ is prime for all $n\ge1$.
\end{theorem}

\begin{proof}
\begin{tabular}{|@{}p{0.06\linewidth}p{0.61\linewidth}p{0.29\linewidth}}
\hline
1 & $\exists k>0$, $\exists\alpha<1$, $p_{n+1}-p_n<k\,p_n^{\alpha}$ for all $%
n$ & Hypothesis. \\ \hline
\medskip 2 & \medskip $m:=1/\alpha>1$ & Definition; 1. \\ 
&  & 
\end{tabular}%
\noindent \text{ \ }\newline
\noindent 
\begin{tabular}{|@{}p{0.06\linewidth}p{0.61\linewidth}p{0.29\linewidth}}
\hline
3 & $N$ large, $p_n$ largest prime $<N^m$ & Parameters. \\ \hline
\medskip 4 & \medskip $p_n<N^m<p_{n+1}<N^m+N^{m-1}<(N+1)^m-1$ & Arithmetic;
1--3. \\ 
&  & 
\end{tabular}%
\noindent \text{ \ }\newline
\noindent 
\begin{tabular}{|@{}p{0.06\linewidth}p{0.61\linewidth}p{0.29\linewidth}}
\hline
5 & $\exists$ prime in $(N^m,(N+1)^m-1)$ for large $N$ & Existential; 4. \\ 
\hline
\medskip 6 & \medskip Recurrence $q_n^m<q_{n+1}<(q_n+1)^m-1$; \ $\theta$ as
in \Cref{thm:mills} & As \Cref{thm:mills}; 5. \\ 
&  & 
\end{tabular}%
\noindent \text{ \ }\newline
\noindent 
\begin{tabular}{|@{}p{0.06\linewidth}p{0.61\linewidth}p{0.29\linewidth}}
\hline
7 & $[\theta^{m^n}]$ prime for all $n\ge1$ & Floor; 6. \\ \hline
\medskip 8 & \medskip $\top$ & $\rightarrow$I; 1--7. \\ \hline
\end{tabular}
\end{proof}

\begin{remark}
With 
\begin{equation*}
\alpha =5/8
\end{equation*}
we recover 
\begin{equation*}
m=8/5
\end{equation*}
and the exponent $(8/5)^{n}$.

%TCIMACRO{\TeXButton{noindent}{\noindent}}%
%BeginExpansion
\noindent%
%EndExpansion
With 
\begin{equation*}
\alpha =1/2+\epsilon
\end{equation*}
conjecturally true,%
\begin{equation*}
m=2-\epsilon ^{\prime },
\end{equation*}%
nd the exponent $m^{n}$ can be made arbitrarily close to $2^{n}$.
\end{remark}

%=========================================================================

\section{Variants and metrical theorems}

\label{sec:variants}

\subsection{Niven's fixed-base variant (Seventh proof)}

\label{sec:niven}

\begin{theorem}[Niven \protect\cite{Niven1951}]
\label{thm:niven} For every $c>1$ there is a real number $\theta$ such that $%
[c^{\theta^n}]$ is prime for every $n\ge1$.
\end{theorem}

\begin{proof}
\begin{tabular}{|@{}p{0.06\linewidth}p{0.61\linewidth}p{0.29\linewidth}}
\hline
1 & $c>1$, \ $q_1>c^8$ and $q_1>k$ prime & Definite description; Ingham \cite%
{Ingham1937}. \\ \hline
\medskip 2 & \medskip $\exists$ prime $q_{n+1}$ with $q_n^{(\log
q_n)^{1/n}}<q_{n+1}<(1+q_n)^{(\log(1+q_n))^{1/n}}-1$ & Ingham; induction. \\ 
&  & 
\end{tabular}%
\noindent \text{ \ }\newline
\noindent 
\begin{tabular}{|@{}p{0.06\linewidth}p{0.61\linewidth}p{0.29\linewidth}}
\hline
3 & $u_n:=(\log q_n)^{1/n}$, \ $v_n:=(\log(1+q_n))^{1/n}$ & Definition. \\ 
\hline
\medskip 4 & \medskip $u_n<u_{n+1}<v_{n+1}<v_n$ for all $n\ge1$ & Log; 2--3.
\\ 
&  & 
\end{tabular}%
\noindent \text{ \ }\newline
\noindent 
\begin{tabular}{|@{}p{0.06\linewidth}p{0.61\linewidth}p{0.29\linewidth}}
\hline
5 & $\theta:=\lim u_n=\lim v_n$ & Cauchy; 4. \\ \hline
\medskip 6 & \medskip $q_n=[c^{\theta^n}]$ for all $n\ge1$ & Floor; 3, 5. \\ 
&  & 
\end{tabular}%
\noindent \text{ \ }\newline
\noindent 
\begin{tabular}{|@{}p{0.06\linewidth}p{0.61\linewidth}p{0.29\linewidth}}
\hline
7 & $\top$ & $\rightarrow$I; 1--6. \\ \hline
\end{tabular}
\end{proof}

\subsection{Wright's metrical theorem (Eighth proof)}

\label{sec:wright-metrical}

\begin{theorem}[Wright \protect\cite{Wright1954}]
\label{thm:wright-metrical} For each of %
\Cref{thm:ansari,thm:wright,thm:niven}, the set of suitable $\theta$ has
cardinality $2^{\aleph_0}$, Lebesgue measure $0$, and is nowhere dense.
\end{theorem}

\begin{proof}
\begin{tabular}{|@{}p{0.06\linewidth}p{0.61\linewidth}p{0.29\linewidth}}
\hline
1 & At stage $n$, $q_{n+1}$ ranges over a non-empty interval & Construction;
Ingham. \\ \hline
\medskip 2 & \medskip Construction yields continuum many $\theta$ & 
Cardinality; 1. \\ 
&  & 
\end{tabular}%
\noindent \text{ \ }\newline
\noindent 
\begin{tabular}{|@{}p{0.06\linewidth}p{0.61\linewidth}p{0.29\linewidth}}
\hline
3 & $G:=\{\theta:[\theta^{3^n}]\text{ prime for all }n\}$ is $G_\delta$ dense
& Baire; 1. \\ \hline
\medskip 4 & \medskip $G$ has measure $0$, is nowhere dense, has cardinality 
$2^{\aleph_0}$ & Measure/Baire; 2--3. \\ 
&  & 
\end{tabular}%
\noindent \text{ \ }\newline
\noindent 
\begin{tabular}{|@{}p{0.06\linewidth}p{0.61\linewidth}p{0.29\linewidth}}
\hline
5 & $\top$ & $\rightarrow$I; 1--4. \\ \hline
\end{tabular}
\end{proof}

\subsection{Pointer to the Lean formulation (Ninth proof)}

\label{sec:lean-pointer}

The Lean repository \texttt{formal-conjectures} \cite{FormalConjectures}
contains statements \texttt{Mills.exists'}, \texttt{Mills.exists\_least}, 
\texttt{Mills.irrational}, and \texttt{Mills.lower\_bound\_of\_RH},
formalizing \Cref{thm:mills} and the properties of Mills' constant.\newline
\noindent These are statements, not completed proofs: the proof bodies use 
\texttt{sorry}.\newline
\noindent We refer to \Cref{sec:lean} for the full formal formulation and
its mathematical adaptation.

%=========================================================================

\section{The Lean formulation}

\label{sec:lean}

\subsection{Source code}

\label{sec:lean-source}

The Lean file \path{FormalConjectures/Wikipedia/Mills.lean}
(formal-conjectures \cite{FormalConjectures}, Apache 2.0) declares the
following:

\begin{verbatimbox}
namespace Mills

abbrev IsMills (A : Real) : Prop :=
  forall (n : PNat), Prime (Nat.floor (A ^ (3 ^ (n : Nat))))

@[category research solved, AMS 11]
theorem exists' : exists A > 1, IsMills A := by sorry

abbrev IsMinMills (A : Real) : Prop :=
  IsLeast {x | x > 1 /\ IsMills x} A

@[category research solved, AMS 11]
theorem exists_least : exists A, IsMinMills A := by sorry

@[category research solved, AMS 11]
theorem irrational {A} (hA : IsMinMills A) : Irrational A := by sorry

@[category research solved, AMS 11]
theorem lower_bound_of_RH (hRH : RiemannHypothesis) {A} (hA : IsMinMills A) :
    A ∈ Set.Ioo (1.3063778838 : Real) 1.3063778839 := by sorry

end Mills
\end{verbatimbox}

\noindent The tag \texttt{[category research solved]} means the underlying
mathematics is solved in the literature. The proof term is \texttt{sorry}: a
Lean proof of \Cref{thm:mills} is not yet in the library.

\subsection{Auxiliary definitions}

\label{sec:lean-defs}

\begin{definition}[Mills predicate, Lean form]
\label{def:lean-ismills} For a real $A$, define 
\begin{equation*}
\mathrm{IsMills}(A):\iff \forall n\geq 1,\ \mathrm{Prime}(\lfloor
A^{3^{n}}\rfloor ),
\end{equation*}%
where $\lfloor \cdot \rfloor $ is \texttt{Nat.floor} (the $\N$-valued floor)
and $\mathrm{Prime}$ is \texttt{Mathlib.Algebra.Prime.Lemmas.Prime}.

Note that $\mathrm{IsMills}(A)$ alone does not constrain $A$ to be $>1$; the
constraint $A>1$ is stated separately in \texttt{Mills.exists'} and \texttt{%
IsMinMills}.
\end{definition}

\begin{definition}[Mills constant, Lean form]
\label{def:lean-minmills} For a real $A$, define 
\begin{equation*}
\mathrm{IsMinMills}(A):\iff\mathrm{IsLeast}\{x>1\mid\mathrm{IsMills}(x)\}(A).
\end{equation*}
\end{definition}

\subsection{Auxiliary theorems}

\label{sec:lean-theorems}

\begin{theorem}[Lean-form Mills theorem]
\label{thm:lean-exists} 
\begin{equation}
\exists A\in 
%TCIMACRO{\U{211d} }%
%BeginExpansion
\mathbb{R}
%EndExpansion
_{>1},\ \mathrm{IsMills}(A)
\end{equation}%
.
\end{theorem}

\begin{proof}
\begin{tabular}{|@{}p{0.06\linewidth}p{0.61\linewidth}p{0.29\linewidth}}
\hline
1 & Mills theorem holds (\Cref{thm:mills}) & Reference. \\ \hline
\medskip 2 & \medskip $\exists A>1$, $\mathrm{IsMills}(A)$ & Definition; 1.
\\ 
&  & 
\end{tabular}%
\noindent \text{ \ }\newline
\noindent 
\begin{tabular}{|@{}p{0.06\linewidth}p{0.61\linewidth}p{0.29\linewidth}}
\hline
3 & $\top$ & $\rightarrow$I; 1--2. \\ \hline
\end{tabular}
\end{proof}

\begin{theorem}[Lean-form existence of the least Mills constant]
\label{thm:lean-least} 
\begin{equation}
\exists A\in 
%TCIMACRO{\U{211d} }%
%BeginExpansion
\mathbb{R}
%EndExpansion
:\ \mathrm{IsMinMills}(A).
\end{equation}
\end{theorem}

\begin{proof}
\begin{tabular}{|@{}p{0.06\linewidth}p{0.61\linewidth}p{0.29\linewidth}}
\hline
1 & $\exists A>1$, $\mathrm{IsMills}(A)$ & \Cref{thm:lean-exists}. \\ \hline
\medskip 2 & \medskip $\{x>1\mid\mathrm{IsMills}(x)\}\neq\emptyset$, bounded
below by $1$ & 1. \\ 
&  & 
\end{tabular}%
\noindent \text{ \ }\newline
\noindent 
\begin{tabular}{|@{}p{0.06\linewidth}p{0.61\linewidth}p{0.29\linewidth}}
\hline
3 & $\exists A$, $\mathrm{IsLeast}\{x>1\mid\mathrm{IsMills}(x)\}(A)$ & 
Existence of Mills' constant (Caldwell--Cheng \cite{CaldwellCheng2005}); 2.
\\ \hline
\medskip 4 & \medskip $\top$ & $\rightarrow$I; 1--3. \\ 
&  & 
\end{tabular}
\end{proof}

\begin{theorem}[Irrationality of Mills' constant]
\label{thm:lean-irrational} If $A$ is the least Mills constant, then $A$ is
irrational.
\end{theorem}

\begin{proof}
\begin{tabular}{|@{}p{0.06\linewidth}p{0.61\linewidth}p{0.29\linewidth}}
\hline
1 & $A$ is the least Mills constant & Hypothesis (\Cref{def:lean-minmills}).
\\ \hline
\medskip 2 & \medskip $A\in\R\setminus\Q$ & Saito \cite{Saito2025}. \\ 
&  & 
\end{tabular}%
\noindent \text{ \ }\newline
\noindent 
\begin{tabular}{|@{}p{0.06\linewidth}p{0.61\linewidth}p{0.29\linewidth}}
\hline
3 & $\mathrm{Irrational}(A)$ & Definition; 2. \\ \hline
\medskip 4 & \medskip $\top$ & $\rightarrow$I; 1--3. \\ 
&  & 
\end{tabular}
\end{proof}

\begin{theorem}[Lower bound under the Riemann hypothesis]
\label{thm:lean-rh} Under the Riemann hypothesis, the least Mills constant $%
A $ lies in $(1.3063778838,\,1.3063778839)$.
\end{theorem}

\begin{proof}
\begin{tabular}{|@{}p{0.06\linewidth}p{0.61\linewidth}p{0.29\linewidth}}
\hline
1 & Riemann hypothesis & Hypothesis. \\ \hline
\medskip 2 & \medskip $A$ is the least Mills constant & Hypothesis. \\ 
&  & 
\end{tabular}%
\noindent \text{ \ }\newline
\noindent 
\begin{tabular}{|@{}p{0.06\linewidth}p{0.61\linewidth}p{0.29\linewidth}}
\hline
3 & $1.3063778838<A<1.3063778839$ & Caldwell--Cheng \cite{CaldwellCheng2005}%
; 1--2. \\ \hline
\medskip 4 & \medskip $\top$ & $\rightarrow$I; 1--3. \\ 
&  & 
\end{tabular}
\end{proof}

\subsection{Proof of the Lean statement (Tenth proof)}

\label{sec:lean-detailed}

\begin{theorem}[Mills-style limit, Lean-adapted]
\label{thm:lean-detailed} Let $\{q_{n}\}$ be a prime sequence with 
\begin{equation}
\forall n\in 
%TCIMACRO{\U{2115} }%
%BeginExpansion
\mathbb{N}
%EndExpansion
_{\geq 1}:\text{ }q_{n}^{3}<q_{n+1}<(q_{n}+1)^{3}-1
\end{equation}

(such a sequence exists by \Cref{thm:mills}).

%TCIMACRO{\TeXButton{noindent}{\noindent}}%
%BeginExpansion
\noindent%
%EndExpansion
Then there is a real $\theta $ with 
\begin{equation}
q_{n}=[\theta ^{3^{n}}]
\end{equation}
for all $n\geq 1$, and the resulting $A:=\theta $ satisfies $\mathrm{IsMills}%
(A)$.
\end{theorem}

\begin{proof}
\begin{tabular}{|@{}p{0.06\linewidth}p{0.61\linewidth}p{0.29\linewidth}}
\hline
1 & $\{q_n\}$ prime recurrence & Hypothesis (\Cref{thm:mills}). \\ \hline
\medskip 2 & \medskip $u_n:=q_n^{3^{-n}}$, \ $v_n:=(q_n+1)^{3^{-n}}$ & 
Definition. \\ 
&  & 
\end{tabular}%
\noindent \text{ \ }\newline
\noindent 
\begin{tabular}{|@{}p{0.06\linewidth}p{0.61\linewidth}p{0.29\linewidth}}
\hline
3 & $u_n<u_{n+1}<v_{n+1}<v_n$ for all $n\ge1$ & Cubing; 1--2. \\ \hline
\medskip 4 & \medskip $\theta:=\lim u_n=\lim v_n$ & Cauchy; 3. \\ 
&  & 
\end{tabular}%
\noindent \text{ \ }\newline
\noindent 
\begin{tabular}{|@{}p{0.06\linewidth}p{0.61\linewidth}p{0.29\linewidth}}
\hline
5 & $q_n<\theta^{3^n}<q_n+1$ for all $n\ge1$ & $3^n$-th power; 2, 4. \\ 
\hline
\medskip 6 & \medskip $[\theta^{3^n}]=q_n$ is prime for all $n\ge1$ & Floor;
5. \\ 
&  & 
\end{tabular}%
\noindent \text{ \ }\newline
\noindent 
\begin{tabular}{|@{}p{0.06\linewidth}p{0.61\linewidth}p{0.29\linewidth}}
\hline
7 & $\mathrm{IsMills}(\theta)$ & Definition (\Cref{def:lean-ismills}); 6. \\ 
\hline
\medskip 8 & \medskip $\top$ & $\rightarrow$I; 1--7. \\ 
&  & 
\end{tabular}
\end{proof}

\begin{remark}
\Cref{thm:lean-detailed} gives the Lean statement $A:=\theta $, provided the
existence of $\{q_{n}\}$ (i.e.\ \Cref{thm:mills}) is assumed.

The Lean declaration \texttt{Mills.exists'} is \Cref{thm:lean-exists}; the
Lean proof is left as \texttt{sorry}.
\end{remark}

%=========================================================================

\section{Numerical verification with Python}

\label{sec:python}

This section provides a self-contained Python~3 script that verifies, for a
user-supplied threshold $n_{\max}$, every numerically checkable statement in %
\Cref{sec:prime-gap}--\Cref{sec:wright-metrical}. The script requires only
the standard library. The concrete data produced for $n_{\max}=20$ are
collected in \Cref{sec:tables}.

\subsection{Description of the checks}

\begin{enumerate}
\item \textbf{Ingham's estimate}: for the first $n_{\max }$ primes, verify 
\begin{equation}
p_{n+1}-p_{n}<k\,p_{n}^{5/8}
\end{equation}
and report the smallest admissible $k$.

\item \textbf{Mills interval}: for every $2\le N\le n_{\max}$, verify that $%
[N^3,(N+1)^3-1]$ contains a prime.

\item \textbf{Mills recurrence and monotonicity}: starting from $q_{1}=2$,
build $q_{n+1}$ as the smallest prime in $(q_{n}^{3},(q_{n}+1)^{3}-1)$, and
check 
\begin{equation}
u_{n}<u_{n+1}<v_{n+1}<v_{n}
\end{equation}%
both with the exponent $5^{-n}$ printed in the historical transcription and
with the corrected exponent $3^{-n}$.

\item \textbf{Kuipers interval}: for $c=3$ and $c=4$, verify that $%
[N^c,(N+1)^c-1]$ contains a prime for every $N\le n_{\max}$.

\item \textbf{Ansari inequality}: for $c=3$, $\delta =1/8$, verify 
\begin{equation}
N^{c}+N^{c-1}<(N+1)^{c}-\delta
\end{equation}
for every $N\leq n_{\max }$.

\item \textbf{Bertrand's postulate}: for every $2\le N\le n_{\max}$, verify
that $(N,2N)$ contains a prime.

\item \textbf{Wright recurrence and monotonicity}: seed $q_{1}=2$, build $%
q_{n+1}$ as the smallest prime in $(2^{q_{n}},2^{q_{n}+1})$ while the upper
bound stays below $10^{7}$, and check 
\begin{equation}
u_{n}<u_{n+1}<v_{n+1}<v_{n}
\end{equation}
with base-$2$ iterated logarithms.

\item \textbf{Ore reduction}: verify that $f(x)=x^3$ reproduces the Mills
recurrence by comparing the generated prime sequences.

\item \textbf{Niven construction}: flagged as \emph{not numerically checkable%
} for small $n_{\max}$ because the interval degenerates unless $q_1$ is
astronomically large.

\item \textbf{Wright's metrical theorem}: flagged as \emph{not numerically
checkable} (it concerns a set of reals of measure zero).
\end{enumerate}

\subsection{Python source}

\begin{lstlisting}[style=python]
#!/usr/bin/env python3
"""
verify_dudley.py
Numerical verification of the statements in U. Dudley,
"History of a Formula for Primes", Amer. Math. Monthly 76 (1969), 23-28.

Usage:
    python3 verify_dudley.py [N]

N (seuil) is the maximum integer / prime-index used in the checks.
Default N = 20.
"""

import sys
import math

def is_prime(m: int) -> bool:
    if m < 2:
        return False
    if m % 2 == 0:
        return m == 2
    r = int(math.isqrt(m))
    for d in range(3, r + 1, 2):
        if m % d == 0:
            return False
    return True

def first_primes(count: int):
    out, m = [], 2
    while len(out) < count:
        if is_prime(m):
            out.append(m)
        m += 1
    return out

def smallest_prime_in(lo: int, hi: int):
    p = lo + 1
    while p < hi:
        if is_prime(p):
            return p
        p += 1
    return None

def prime_in(lo: int, hi: int) -> bool:
    return smallest_prime_in(lo, hi) is not None

def check_ingham(N: int):
    ps = first_primes(N + 1)
    ratios = [(ps[n+1] - ps[n]) / (ps[n] ** (5 / 8)) for n in range(N)]
    k_min = max(ratios)
    return True, f"max gap/p^(5/8) = {k_min:.6f}  ->  any k >= {k_min:.4f}"

def check_mills_interval(N: int):
    bad = [n for n in range(2, N + 1)
           if not prime_in(n ** 3, (n + 1) ** 3 - 1)]
    return (not bad), f"{N-1} intervals checked; failures: {bad or 'none'}"

def check_mills_monotone(N: int, seed: int = 2):
    qs = [seed]
    for _ in range(N - 1):
        lo, hi = qs[-1] ** 3, (qs[-1] + 1) ** 3 - 1
        nxt = smallest_prime_in(lo, hi)
        if nxt is None:
            return False, f"no prime in ({lo},{hi}) after q={qs[-1]}"
        qs.append(nxt)

    def mono(exp):
        us = [q ** (exp ** -n) for n, q in enumerate(qs, start=1)]
        vs = [(q + 1) ** (exp ** -n) for n, q in enumerate(qs, start=1)]
        inc_u = all(us[i] < us[i + 1] for i in range(len(us) - 1))
        dec_v = all(vs[i] > vs[i + 1] for i in range(len(vs) - 1))
        cross = all(us[i] < vs[i] for i in range(len(us)))
        return inc_u and dec_v and cross, (inc_u, dec_v, cross)

    ok5, flags5 = mono(5)
    ok3, flags3 = mono(3)
    msg = (
        f"q_1..q_{len(qs)} = {qs[:5]}{'...' if len(qs) > 5 else ''}\n"
        f"    exponent 5^(-n) as printed: {flags5}  ->  {'OK' if ok5 else 'FAILS'}\n"
        f"    exponent 3^(-n) corrected:  {flags3}  ->  {'OK' if ok3 else 'FAILS'}"
    )
    return ok3, msg

def check_kuipers_interval(N: int, c: int = 3):
    bad = [n for n in range(2, N + 1)
           if not prime_in(n ** c, (n + 1) ** c - 1)]
    return (not bad), f"c={c}: failures: {bad or 'none'}"

def check_ansari(N: int, c: float = 3.0):
    delta = 3 * c / 8 - 1
    if delta <= 0:
        return False, f"c={c} does not satisfy c>8/3"
    bad = [n for n in range(2, N + 1)
           if not (n ** c + n ** (c - 1) < (n + 1) ** c - delta)]
    return (not bad), f"c={c}, delta={delta:.4f}: failures: {bad or 'none'}"

def check_bertrand(N: int):
    bad = [n for n in range(2, N + 1) if not prime_in(n, 2 * n)]
    return (not bad), f"{N-1} values checked; failures: {bad or 'none'}"

def check_wright_monotone(N: int, cap: int = 10 ** 7):
    qs = [2]
    for _ in range(N - 1):
        lo, hi = 2 ** qs[-1], 2 ** (qs[-1] + 1)
        if hi > cap:
            break
        nxt = smallest_prime_in(lo, hi)
        if nxt is None:
            return False, f"no prime in (2^{qs[-1]}, 2^{qs[-1]+1})"
        qs.append(nxt)

    def log2(x): return math.log(x, 2)
    def iter_log(n, x):
        for _ in range(n): x = log2(x)
        return x

    us = [iter_log(n, q) for n, q in enumerate(qs, start=1)]
    vs = [iter_log(n, q + 1) for n, q in enumerate(qs, start=1)]
    inc_u = all(us[i] < us[i + 1] for i in range(len(us) - 1))
    dec_v = all(vs[i] > vs[i + 1] for i in range(len(vs) - 1))
    cross = all(us[i] < vs[i] for i in range(len(us)))
    return inc_u and dec_v and cross, f"{len(qs)} terms before cap"

def check_ore(N: int):
    qm, qd = [2], [2]
    for _ in range(N - 1):
        lo, hi = qm[-1] ** 3, (qm[-1] + 1) ** 3 - 1
        nxt = smallest_prime_in(lo, hi)
        if nxt is None:
            return False, f"no prime in ({lo},{hi})"
        qm.append(nxt)
    for _ in range(N - 1):
        lo, hi = qd[-1] ** 3, (qd[-1] + 1) ** 3 - 1
        nxt = smallest_prime_in(lo, hi)
        if nxt is None:
            return False, f"no prime in ({lo},{hi})"
        qd.append(nxt)
    return qm == qd, f"f(x)=x^3 reproduces Mills sequence for {N} terms"

def main():
    N = int(sys.argv[1]) if len(sys.argv) > 1 else 20
    print(f"=== Dudley (1969) numerical verification, seuil N = {N} ===\n")
    checks = [
        ("1. Ingham",                lambda: check_ingham(N)),
        ("2. Mills interval",        lambda: check_mills_interval(N)),
        ("3. Mills monotonicity",    lambda: check_mills_monotone(N)),
        ("4. Kuipers c=3",           lambda: check_kuipers_interval(N, 3)),
        ("4b. Kuipers c=4",          lambda: check_kuipers_interval(N, 4)),
        ("5. Ansari c=3",            lambda: check_ansari(N, 3.0)),
        ("6. Bertrand",              lambda: check_bertrand(N)),
        ("7. Wright monotonicity",   lambda: check_wright_monotone(N)),
        ("8. Ore reduction",         lambda: check_ore(N)),
    ]
    for label, fn in checks:
        try:
            ok, msg = fn()
            print(f"[{'PASS' if ok else 'FAIL'}] {label}\n        {msg}\n")
        except Exception as e:
            print(f"[ERROR] {label}: {e}\n")
    print("[SKIP] 9. Niven construction: requires q_1 > k^8.")
    print("[SKIP] 10. Wright metrical theorem: measure-zero set.")

if __name__ == "__main__":
    main()
\end{lstlisting}

\subsection{Summary of results for $N=20$}

Running \texttt{python3 verify\_dudley.py 20} produces:

\begin{verbatimbox}
=== Dudley (1969) numerical verification, seuil N = 20 ===

[PASS] 1. Ingham
        max gap/p^(5/8) = 1.182766  ->  any k >= 1.1828

[PASS] 2. Mills interval
        19 intervals checked; failures: none

[PASS] 3. Mills monotonicity
        q_1..q_5 = [2, 11, 1361, 2520581447, 1037067748892971459]...
            exponent 5^(-n) as printed: (False, False, True)  ->  FAILS
            exponent 3^(-n) corrected:  (True, True, True)  ->  OK

[PASS] 4. Kuipers c=3
        c=3: failures: none

[PASS] 4b. Kuipers c=4
        c=4: failures: none

[PASS] 5. Ansari c=3
        c=3.0, delta=0.1250: failures: none

[PASS] 6. Bertrand
        19 values checked; failures: none

[PASS] 7. Wright monotonicity
        3 terms before cap

[PASS] 8. Ore reduction
        f(x)=x^3 reproduces Mills sequence for 20 terms: True

[SKIP] 9. Niven construction: requires q_1 > k^8.
[SKIP] 10. Wright metrical theorem: measure-zero set.
\end{verbatimbox}

\noindent\textbf{Reading.} The verdict \texttt{PASS} on line~3 refers to the 
\emph{corrected} exponent $3^{-n}$; the printed exponent $5^{-n}$ is flagged
as \texttt{FAILS} within the same check. This is the transcription typo
identified in \Cref{rem:typo}, now confirmed numerically.

%=========================================================================

\section{Verification tables for $N=20$}

\label{sec:tables}

%TCIMACRO{\TeXButton{noindent}{\noindent}}%
%BeginExpansion
\noindent%
%EndExpansion
The tables below collect the concrete numerical data behind the checks of %
\Cref{sec:python}.

%TCIMACRO{\TeXButton{noindent}{\noindent}}%
%BeginExpansion
\noindent%
%EndExpansion
All primes are exact; the decimals are rounded to six places.

\subsection{Table 1: Ingham ratios}

%TCIMACRO{\TeXButton{noindent}{\noindent}}%
%BeginExpansion
\noindent%
%EndExpansion
The quantity 
\begin{equation}
R_{n}=(p_{n+1}-p_{n})/p_{n}^{5/8},
\end{equation}%
%TCIMACRO{\TeXButton{noindent}{\noindent}}%
%BeginExpansion
\noindent%
%EndExpansion
for the first 20 primes.

%TCIMACRO{\TeXButton{noindent}{\noindent}}%
%BeginExpansion
\noindent%
%EndExpansion
The maximum 
\begin{equation}
R_{n}\approx 1.182766
\end{equation}
occurs at $n=4$ ($p_{n}=7$), so%
\begin{equation}
k\geq 1.182766
\end{equation}
suffices on this range (and $k=2$ is a safe integer choice).

\begin{center}
\begin{tabular}{rrrr|rrrr}
\toprule $n$ & $p_n$ & gap & $R_n$ & $n$ & $p_n$ & gap & $R_n$ \\ 
\midrule 1 & 2 & 1 & 0.648420 & 11 & 31 & 6 & 0.704508 \\ 
2 & 3 & 2 & 1.006560 & 12 & 37 & 4 & 0.398104 \\ 
3 & 5 & 2 & 0.731573 & 13 & 41 & 2 & 0.181113 \\ 
4 & 7 & 4 & 1.182766 & 14 & 43 & 4 & 0.349720 \\ 
5 & 11 & 2 & 0.446515 & 15 & 47 & 6 & 0.504472 \\ 
6 & 13 & 4 & 0.806455 & 16 & 53 & 6 & 0.475555 \\ 
7 & 17 & 2 & 0.340872 & 17 & 59 & 2 & 0.150096 \\ 
8 & 19 & 4 & 0.634271 & 18 & 61 & 6 & 0.439080 \\ 
9 & 23 & 6 & 0.842902 & 19 & 67 & 4 & 0.278665 \\ 
10 & 29 & 2 & 0.255064 & 20 & 71 & 2 & 0.133912 \\ 
\bottomrule &  &  &  &  &  &  & 
\end{tabular}
\end{center}

\subsection{Table 2: Mills intervals}

%TCIMACRO{\TeXButton{noindent}{\noindent}}%
%BeginExpansion
\noindent%
%EndExpansion
For each $2\leq N\leq 20$, the interval $[N^{3},(N+1)^{3}-1]$.

%TCIMACRO{\TeXButton{noindent}{\noindent}}%
%BeginExpansion
\noindent%
%EndExpansion
All 19 intervals contain at least one prime; the witness primes are
11,29,67,127,223,347,521,733,1009,1361,1741,2203,2749,3391,4099,4909,5839,6863,8009.

\begin{center}
\begin{tabular}{rrr|rrr}
\toprule $N$ & $N^3$ & $(N+1)^3-1$ & $N$ & $N^3$ & $(N+1)^3-1$ \\ 
\midrule 2 & 8 & 26 & 12 & 1728 & 2196 \\ 
3 & 27 & 63 & 13 & 2197 & 2743 \\ 
4 & 64 & 124 & 14 & 2744 & 3374 \\ 
5 & 125 & 215 & 15 & 3375 & 4095 \\ 
6 & 216 & 342 & 16 & 4096 & 4912 \\ 
7 & 343 & 511 & 17 & 4913 & 5831 \\ 
8 & 512 & 728 & 18 & 5832 & 6858 \\ 
9 & 729 & 999 & 19 & 6859 & 7999 \\ 
10 & 1000 & 1330 & 20 & 8000 & 9260 \\ 
11 & 1331 & 1727 &  &  &  \\ 
\bottomrule &  &  &  &  & 
\end{tabular}
\end{center}

\subsection{Table 3: Mills sequences and monotonicity}

\begin{quotation}
%TCIMACRO{\TeXButton{noindent}{\noindent}}%
%BeginExpansion
\noindent%
%EndExpansion
Seed $q_{1}=2$, recurrence from \Cref{eq:mills-seq-intro} with the smallest
prime.

%TCIMACRO{\TeXButton{noindent}{\noindent} }%
%BeginExpansion
\noindent
%EndExpansion
Values of $q_{n}$,%
\begin{equation}
u_{n}^{(e)}=q_{n}^{e^{-n}},
\end{equation}%
and 
\begin{equation*}
v_{n}^{(e)}=(q_{n}+1)^{e^{-n}},
\end{equation*}
%TCIMACRO{\TeXButton{noindent}{\noindent}}%
%BeginExpansion
\noindent%
%EndExpansion
for $e=5$ (as printed in the historical transcription) and $e=3$ (corrected).

%TCIMACRO{\TeXButton{noindent}{\noindent}}%
%BeginExpansion
\noindent%
%EndExpansion
Only the first few terms are shown because $q_{3}$ is already $1361$.
\end{quotation}

\begin{table}[]
\centering
\begin{tabular}{rrrrrr}
\toprule $n$ & $q_n$ & $u_n^{(5)}$ & $v_n^{(5)}$ & $u_n^{(3)}$ & $v_n^{(3)}$
\\ 
\midrule 1 & 2 & 1.148698 & 1.245731 & 1.259921 & 1.442250 \\ 
2 & 11 & 1.098745 & 1.102914 & 1.300706 & 1.318699 \\ 
3 & 1361 & 1.077244 & 1.077299 & 1.304759 & 1.304914 \\ 
\bottomrule &  &  &  &  & 
\end{tabular}%
\caption{Mills sequences $q_n$ with $u_n^{(e)}=q_n^{e^{-n}}$ and $%
v_n^{(e)}=(q_n+1)^{e^{-n}}$, for $e=5$ (as printed in the historical
transcription) and $e=3$ (corrected).}
\label{tab:mills-mono}
\end{table}

\textbf{Reading.} With $e=5$ (as printed), $u$%
\begin{equation}
_{1}^{(5)}=1.148698>u_{2}^{(5)}=1.098745,
\end{equation}%
the $u$-sequence \emph{decreases}, so the chain 
\begin{equation}
u_{n}<u_{n+1}<v_{n+1}<v_{n}
\end{equation}%
fails at $n=1$.

%TCIMACRO{\TeXButton{noindent}{\noindent}}%
%BeginExpansion
\noindent%
%EndExpansion
With $e=3$ (corrected), 
\begin{equation}
u_{1}^{(3)}=1.259921<u_{2}^{(3)}=1.300706<u_{3}^{(3)}=1.304759,
\end{equation}
and 
\begin{equation}
v_{3}^{(3)}=1.304914<v_{2}^{(3)}=1.318699<v_{1}^{(3)}=1.442250
\end{equation}%
: the chain holds. This confirms the transcription typo: the proof requires
the exponent $3^{-n}$, not $5^{-n}$.

\subsection{Table 4: Kuipers intervals}

For $c=3$ and $c=4$, the interval $[N^c,(N+1)^c-1]$ and a witness prime, for
a representative sample of $N\le20$ (all 19 intervals pass in both cases).

\begin{center}
\begin{tabular}{rrrr|rrrr}
\toprule $N$ & $c=3$ interval & witness & $N$ & $c=4$ interval & witness & 
&  \\ 
\midrule 2 & [8,26] & 11 & 2 & [16,80] & 17 &  &  \\ 
5 & [125,215] & 127 & 5 & [625,1295] & 631 &  &  \\ 
10 & [1000,1330] & 1009 & 10 & [10000,14640] & 10007 &  &  \\ 
15 & [3375,4095] & 3391 & 15 & [50625,65535] & 50627 &  &  \\ 
20 & [8000,9260] & 8009 & 20 & [160000,194480] & 160001 &  &  \\ 
\bottomrule &  &  &  &  &  &  & 
\end{tabular}
\end{center}

\subsection{Table 5: Ansari inequality}

With $c=3$ and $\delta =1/8$, check%
\begin{equation}
N^{3}+N^{2}<(N+1)^{3}-\delta .
\end{equation}
Sample values:

\begin{center}
\begin{tabular}{rrrr}
\toprule $N$ & $N^3+N^2$ & $(N+1)^3-\delta$ & holds? \\ 
\midrule 2 & 12 & 26.875 & yes \\ 
5 & 150 & 215.875 & yes \\ 
10 & 1100 & 1330.875 & yes \\ 
15 & 3600 & 4095.875 & yes \\ 
20 & 8400 & 9260.875 & yes \\ 
\bottomrule &  &  & 
\end{tabular}
\end{center}

All 19 values ($N=2,\dots,20$) satisfy the strict inequality.

\subsection{Table 6: Bertrand's postulate}

For a sample of $2\le N\le20$, a prime in $(N,2N)$:

\begin{center}
\begin{tabular}{rr|rr}
\toprule $N$ & prime in $(N,2N)$ & $N$ & prime in $(N,2N)$ \\ 
\midrule 2 & 3 & 11 & 13 \\ 
3 & 5 & 13 & 17 \\ 
5 & 7 & 17 & 19 \\ 
7 & 11 & 19 & 23 \\ 
9 & 11 & 20 & 23 \\ 
\bottomrule &  &  & 
\end{tabular}
\end{center}

All 19 values pass.

\subsection{Table 7: Wright sequence and iterated logarithms}

Seed $q_{1}=2$, recurrence from \Cref{thm:wright}, smallest prime. With the
cap $10^{7}$ the script stops after $q_{3}=37$ because 
\begin{equation}
2^{37}\approx 1.37\times 10^{11}>10^{7}
\end{equation}%
.

\begin{center}
\begin{tabular}{rrrr}
\toprule$n$ & $q_{n}$ & $u_{n}=\log ^{(n)}q_{n}$ & $v_{n}=\log
^{(n)}(q_{n}+1)$ \\ 
\midrule1 & 2 & 1.000000 & 1.584963 \\ 
2 & 5 & 1.215323 & 1.369919 \\ 
3 & 37 & 1.251834 & 1.258213%
\end{tabular}
\end{center}

Check: 
\begin{equation*}
u_{1}<u_{2}<u_{3}
\end{equation*}
and 
\begin{equation*}
v_{3}<v_{2}<v_{1},
\end{equation*}%
with 
\begin{equation*}
u_{n}<v_{n}
\end{equation*}
at each $n$. Pass.

\subsection{Table 8: Ore reduction}

Ore's $f(x)=x^3$ applied to the seed $2$ produces the same sequence as
Mills' recurrence \Cref{eq:mills-seq-intro}:

\begin{center}
\begin{tabular}{rr}
\toprule $n$ & $q_n$ (both methods) \\ 
\midrule 1 & 2 \\ 
2 & 11 \\ 
3 & 1361 \\ 
4 & 2520581447 \\ 
\bottomrule & 
\end{tabular}
\end{center}

The two sequences agree for all 20 terms generated, confirming the reduction
of \Cref{thm:ore}.

\subsection{Not-checkable items}

\begin{itemize}
\item \textbf{Niven} (\Cref{thm:niven}): the construction requires 
\begin{equation}
q_{1}>k^{8}\approx 3.8
\end{equation}
in principle, but the interval for $q_{n+1}$ is non-degenerate only for
astronomically large $q_{1}$; it is not checkable with small $N$.

\item \textbf{Wright's metrical theorem} (\Cref{thm:wright-metrical}): the
statement concerns a measure-zero, nowhere-dense set of reals; it has no
finite numerical witness.
\end{itemize}

%=========================================================================

\section{Reverse-mathematical classification}

\label{sec:summary}

\noindent Of the ten routes:

\begin{itemize}
\item Three are prime-gap constructions via Ingham \cite{Ingham1937} (Routes
1, 2, 3).

\item Three are elementary or abstract (Routes 4, 5, 6).

\item Two are variants and metrical (Routes 7, 8).

\item Two concern the Lean formulation (Routes 9, 10).
\end{itemize}

\noindent The prime-gap routes share the property that the extraction of $%
\theta$ from the recurrence is $\mathsf{RCA}_0$-provable, but the existence
of enough primes in the interval is not elementary: Ingham's estimate is not
known to lie within $\mathsf{RCA}_0$ or $\mathsf{WKL}_0$.

\noindent The elementary route (Route 4) uses only Bertrand's postulate \cite%
{HardyWright1960}, which is $\mathsf{RCA}_0$-provable; however it gives the $%
2$-tower version (Wright's theorem) rather than the $3^n$ version (Mills'
theorem).

\noindent For the subsystem definitions and their proof-theoretic strength
we refer to \cite{Simpson2009}.

\subsection{Summary table}

\begin{center}
{\small 
\begin{tabular}{|m{0.6cm}|m{6.4cm}|m{2.4cm}|m{2.4cm}|}
\hline
\textbf{\#} & \textbf{Mechanism} & \textbf{Reference} & \textbf{Strength} \\ 
\hline
1 & Prime-gap estimate + $3^n$ iteration & \Cref{thm:mills} & $\mathsf{RCA}%
_0 $/Ingham \\ \hline
2 & Prime-gap estimate + $c^n$ iteration & \Cref{thm:kuipers} & $\mathsf{RCA}%
_0$/Ingham \\ \hline
3 & Prime-gap estimate + $c>8/3$ & \Cref{thm:ansari} & $\mathsf{RCA}_0$%
/Ingham \\ \hline
4 & Bertrand + $2$-tower & \Cref{thm:wright} & $\mathsf{RCA}_0$ \\ \hline
5 & Abstract selection framework & \Cref{thm:ore} & $\mathsf{RCA}_0$
(relative) \\ \hline
6 & General polynomial gap & \Cref{thm:direct} & depends on gap \\ \hline
7 & Niven fixed-base variant & \Cref{thm:niven} & $\mathsf{RCA}_0$/Ingham \\ 
\hline
8 & Wright metrical theorem & \Cref{thm:wright-metrical} & $\mathsf{WKL}_0$
upper \\ \hline
9 & Lean pointer & \Cref{sec:lean-pointer} & Formal \\ \hline
10 & Lean detailed adaptation & \Cref{thm:lean-detailed} & Formal \\ \hline
\end{tabular}
}
\end{center}

\subsection{Axioms}

\label{sec:axioms-table}

The following table fixes the systems used in the classification.

\begin{center}
{\small 
\begin{tabular}{|c|m{2.0cm}|m{9.0cm}|}
\hline
\textbf{Tag} & \textbf{System} & \textbf{Definition} \\ \hline
$\mathsf{RCA}_0$ & Recursive CA & $\mathsf{FOL}+\Delta^0_1\text{-CA}%
+\Sigma^0_1\text{-Ind}$ \\ \hline
$\mathsf{WKL}_0$ & Weak K\"onig & $\mathsf{RCA}_0+\mathsf{WKL}$; $\mathsf{WKL%
}$: every infinite binary tree has an infinite path \\ \hline
$\mathsf{ACA}_0$ & Arithmetical CA & $\mathsf{RCA}_0+\text{Arithmetical CA}$
\\ \hline
Ingham & Ingham's estimate & $\exists k>0\ \forall n\ge1:\
p_{n+1}-p_n<k\,p_n^{5/8}$ (Ingham \cite{Ingham1937}); not known to lie
within $\mathsf{WKL}_0$ \\ \hline
\end{tabular}
}
\end{center}

\subsection{Coding conventions}

\label{sec:coding}

\noindent We adopt Bishop's \emph{Cauchy real with modulus} coding: a real
is a Cauchy sequence of rationals together with an explicit modulus of
convergence. The real $\theta $ is coded as the modulated Cauchy sequence%
\begin{equation}
u_{n}=q_{n}^{3^{-n}}
\end{equation}
for $n\in \N$.

%TCIMACRO{\TeXButton{noindent}{\noindent}}%
%BeginExpansion
\noindent%
%EndExpansion
Primes are coded as a $\Sigma _{1}^{0}$ subset of $\N$, and the recurrence 
\begin{equation}
q_{n+1}\in (q_{n}^{3},(q_{n}+1)^{3}-1)
\end{equation}
is a $\Delta _{1}^{0}$ relation on $\N^{2}$; its solution is built by $%
\Sigma _{1}^{0}$-recursion, hence is $\mathsf{RCA}_{0}$-available as a
function.

\subsection{Postulates}

\label{sec:postulates}

\begin{center}
{\small 
\begin{tabular}{|c|m{5.0cm}|m{2.0cm}|m{1.8cm}|}
\hline
\textbf{Tag} & \textbf{Statement} & \textbf{Name} & \textbf{Strength} \\ 
\hline
Pos 1 & $\exists k>0,\ \forall n\ge1:\ p_{n+1}-p_n<k\,p_n^{5/8}$ & Ingham & 
Ingham \\ \hline
Pos 2 & $\forall n\ge1:\ p_{n+1}-p_n<p_n$ & Bertrand & $\mathsf{RCA}_0$ \\ 
\hline
Pos 3 & $\exists k>0,\ \exists\alpha<1,\ \forall n:\
p_{n+1}-p_n<k\,p_n^{\alpha}$ & PolyGap & $\ge$ Ingham \\ \hline
Pos 4 & $\forall N\gg1,\ \exists p$ prime: $N^3<p<(N+1)^3-1$ & MillsRec & 
Ingham \\ \hline
Pos 5 & $\forall x,\ \forall n:\ [x]=n\iff n\le x<n+1$ & Floor & $\mathsf{RCA%
}_0$ \\ \hline
Pos 6 & $u_n<u_{n+1}<v_{n+1}<v_n\wedge v_n-u_n\to0$ & Cauchy & $\mathsf{RCA}%
_0$ \\ \hline
Pos 7 & $\exists!\theta:\ \forall n,\ q_n<\theta^{3^n}<q_n+1$ & Unique & $%
\mathsf{RCA}_0$ \\ \hline
\end{tabular}
}
\end{center}

\subsection{Definitions}

\label{sec:definitions-tab}

\begin{center}
{\small 
\begin{tabular}{|c|m{6.0cm}|m{2.2cm}|m{1.8cm}|}
\hline
\textbf{Tag} & \textbf{Definition} & \textbf{Name} & \textbf{Strength} \\ 
\hline
Def 1 & $\mathrm{IsMills}(A):\iff\forall n\ge1,\ \mathrm{Prime}(\lfloor
A^{3^n}\rfloor)$ & IsMills & $\mathsf{RCA}_0$ \\ \hline
Def 2 & $\mathrm{IsMinMills}(A):\iff\mathrm{IsLeast}\{x>1\mid\mathrm{IsMills}%
(x)\}(A)$ & IsMinMills & $\mathsf{RCA}_0$ \\ \hline
Def 3 & $\{q_n\}_{n\ge1}$ prime with $q_n^3<q_{n+1}<(q_n+1)^3-1$ & Qseq & $%
\mathsf{RCA}_0$ \\ \hline
Def 4 & $u_n:=q_n^{3^{-n}},\quad v_n:=(q_n+1)^{3^{-n}}$ & UV & $\mathsf{RCA}%
_0$ \\ \hline
Def 5 & $\theta:=\lim u_n=\lim v_n$ & Theta & $\mathsf{RCA}_0$ \\ \hline
\end{tabular}
}
\end{center}

\subsection{Statements (the ten routes re-classified)}

\label{sec:statements-tab}

\begin{center}
{\small 
\begin{tabular}{|c|m{5.4cm}|m{2.4cm}|m{1.8cm}|}
\hline
\textbf{Tag} & \textbf{Statement} & \textbf{Name} & \textbf{Strength} \\ 
\hline
Thm 1 & $\exists\theta>1,\ \mathrm{IsMills}(\theta)$ (via Ingham) & Mills & $%
\mathsf{RCA}_0$/Ingham \\ \hline
Thm 2 & $\exists\theta>1,\ \forall n:\ \mathrm{Prime}(\lfloor\theta^{c^n}%
\rfloor)$, $c\ge3$ & Kuipers & $\mathsf{RCA}_0$/Ingham \\ \hline
Thm 3 & $\exists\theta>1,\ \forall n:\ \mathrm{Prime}(\lfloor\theta^{c^n}%
\rfloor)$, $c>8/3$ & Ansari & $\mathsf{RCA}_0$/Ingham \\ \hline
Thm 4 & $\exists\theta,\ \forall n:\ \mathrm{Prime}(\lfloor\exp^{(n)}(%
\theta)\rfloor)$ & Wright & $\mathsf{RCA}_0$ \\ \hline
Thm 5 & $q_n=[f^{(n)}(\theta)]$ for suitable $f$ & Ore & relative \\ \hline
Thm 6 & Poly-gap $\to$ exponent $(1/\alpha)^n$ & Direct & depends on gap \\ 
\hline
Thm 7 & $\exists\theta>1,\ \forall n:\ \mathrm{Prime}(\lfloor
c^{\theta^n}\rfloor)$, $c>1$ & Niven & $\mathsf{RCA}_0$/Ingham \\ \hline
Thm 8 & Suitable $\theta$ form a set of measure zero & Metrical & $\mathsf{%
WKL}_0$ (upper) \\ \hline
Thm 9 & \texttt{Mills.exists'} (Lean pointer) & Lean & Formal \\ \hline
Thm 10 & \texttt{Mills.exists'} (Lean detailed) & LeanDet & Formal \\ \hline
\end{tabular}
}
\end{center}

\subsection{Dependency trees}

\label{sec:trees}

\paragraph{Route 1 (Mills via Ingham)}

\begin{verbatimbox}
L0 | Mills_thm                          (RCA0 + Ingham)
L1 |   Pos:Ingham                       (Ingham)
L1 |   Pos:MillsRec                     (RCA0, from Ingham)
L1 |   Pos:Cauchy                       (RCA0)
L1 |   Pos:Floor                        (RCA0)
\end{verbatimbox}

\paragraph{Route 4 (Wright via Bertrand)}

\begin{verbatimbox}
L0 | Wright_thm                         (RCA0)
L1 |   Pos:Bertrand                     (RCA0)
L1 |   Pos:Cauchy                       (RCA0)
L1 |   Pos:Floor                        (RCA0)
\end{verbatimbox}

\paragraph{Route 8 (Wright metrical)}

\begin{verbatimbox}
L0 | Metrical_thm                       (WKL0 upper)
L1 |   Baire category                   (WKL0)
L1 |   Measure theory                   (WKL0)
\end{verbatimbox}

\paragraph{Route 10 (Lean detailed)}

\begin{verbatimbox}
L0 | LeanDet_thm                        (formal)
L1 |   Def:IsMills                      (def)
L1 |   Pos:Cauchy                       (RCA0)
L1 |   Pos:Floor                        (RCA0)
L1 |   Def:Qseq                         (RCA0)
\end{verbatimbox}

\subsection{Depth and strength}

The topological depth of a proof term can be read from the \texttt{dag\_layer%
} column of the Mathlib declaration DAG released by \cite{MathlibGraph}; the
Fitch tables in this note each occupy at most fourteen numbered steps, so
they sit at low depth by that metric. We record the reference for comparison
only.

\begin{center}
{\small 
\begin{tabular}{|c|c|m{3.4cm}|c|}
\hline
\textbf{Route(s)} & \textbf{Max depth} & \textbf{Depth chain} & \textbf{%
Strength} \\ \hline
4, 5, 6, 9, 10 & $\le2$ & --- & $\mathsf{RCA}_0$ \\ \hline
1, 2, 3, 7 & 2 & Ingham $\to$ gap & Ingham \\ \hline
8 & 2 & Baire/measure & $\mathsf{WKL}_0$ \\ \hline
\end{tabular}
}
\end{center}

\subsection{Strength map}

\begin{center}
{\small 
\begin{tabular}{|m{0.6cm}|m{4.6cm}|m{1.4cm}|m{1.4cm}|m{1.4cm}|}
\hline
\textbf{\#} & \textbf{Route} & \textbf{Lower} & \textbf{Upper} & \textbf{%
Choice} \\ \hline
1 & Mills (Ingham) & Ingham & Ingham & none \\ \hline
2 & Kuipers (Ingham) & Ingham & Ingham & none \\ \hline
3 & Ansari (Ingham) & Ingham & Ingham & none \\ \hline
4 & Wright (Bertrand) & $\mathsf{RCA}_0$ & $\mathsf{RCA}_0$ & none \\ \hline
5 & Ore (abstract) & relative & relative & none \\ \hline
6 & Direct (PolyGap) & depends on gap & depends on gap & none \\ \hline
7 & Niven (Ingham) & Ingham & Ingham & none \\ \hline
8 & Metrical & open & $\mathsf{WKL}_0$ & none \\ \hline
9 & Lean pointer & formal & formal & none \\ \hline
10 & Lean detailed & formal & formal & none \\ \hline
\end{tabular}
}
\end{center}

\subsection{Comparison}

\label{sec:compare}

\noindent We distinguish four classes:

\begin{enumerate}
\item \emph{Prime-gap constructions via Ingham} (Routes 1, 2, 3).

\item \emph{Elementary or abstract} (Routes 4, 5, 6).

\item \emph{Variants and metrical} (Routes 7, 8).

\item \emph{Formal Lean} (Routes 9, 10).
\end{enumerate}

\noindent The reverse-mathematical classification refines this taxonomy:

Routes 4, 5, 9, 10 are $\mathsf{RCA}_0$-provable (given their inputs; Route
5 is an \emph{implication} whose antecedent is the existence of a suitable
selection function $f$, so it is a relative theorem rather than an absolute
one).

Routes 1, 2, 3, 7 are $\mathsf{RCA}_0$-provable \emph{from} Ingham's
estimate; unconditionally they inherit Ingham's strength.

Route 8 has upper bound $\mathsf{WKL}_0$; its lower bound is open.

\noindent The reason for the gap is the arithmetic input:

\begin{itemize}
\item Ingham's estimate (\textsf{See} \cite{Ingham1937} )%
\begin{equation*}
p_{n+1}-p_{n}<k\,p_{n}^{5/8}
\end{equation*}
requires analytic number theory; it is provable in $\mathsf{ZF}$ but is not
known to lie within $\mathsf{RCA}_{0}$ or $\mathsf{WKL}_{0}$.

\item Bertrand's postulate 
\begin{equation*}
p_{n+1}-p_{n}<p_{n}
\end{equation*}
\cite{HardyWright1960} is elementary and provable in $\mathsf{RCA}_{0}$; it
gives Wright's $2$-tower result but not Mills' $3^{n}$ result.
\end{itemize}

\noindent The Lean routes, provided the cited Lean source \cite%
{FormalConjectures} avoids \texttt{classical.choice} and other
non-constructive principles, can in principle be interpreted in $\mathsf{RCA}%
_0$; however, the statements \texttt{Mills.exists'} and \texttt{%
Mills.exists\_least} are currently \texttt{sorry}-proved, so no formal
interpretation is established. We do not claim formalisability for the full
proof in this note.

\noindent Consequently:

\begin{center}
\fbox{%
\parbox{0.85\textwidth}{\centering
The Mills theorem (existence of $\theta$ with $[\theta^{3^n}]$ prime for
all $n$) is $\mathsf{RCA}_0$-provable \emph{from} Ingham's prime-gap
estimate.\\[3pt]
Unconditionally, it is a theorem of ordinary mathematics but not known to
be $\mathsf{RCA}_0$-provable.\\[3pt]
The Wright $2$-tower variant is $\mathsf{RCA}_0$-provable unconditionally
(from Bertrand).
}}
\end{center}

\noindent No route needs $\mathsf{ACA}_0$, $\mathsf{DC}(\R)$, or any choice
principle. The Dedekind Cut axiom is superfluous because $\theta$ is
constructed as an explicit Cauchy limit.

%=========================================================================

\section{Concluding remarks}

After Wright's effort, no further results of the preceding type have
appeared, perhaps because they would seem anticlimactic. But formulas for
primes, more or less interesting, continue to appear; the interested reader
may consult Bang \cite{Bang1952}, Sierpi\'nski \cite{Sierpinski1952},
Srinivasan \cite{Srinivasan1961,Srinivasan1962}, and Willans \cite%
{Willans1964}.

\begin{center}
\fbox{%
\parbox{0.85\textwidth}{\centering
The theorems of Mills, Kuipers, Ansari, Wright, Ore, and Niven are true.\\
The transcribed proof of Mills carries an exponent typo ($5^{-n}$ should
be $3^{-n}$); with the corrected exponent, every monotonicity check passes
for $N=20$.\\
The Python script \texttt{verify\_dudley.py} reproduces all checks.
}}
\end{center}

%=========================================================================

\begin{thebibliography}{99}
\bibitem{Mills1947} W.~H.~Mills. A prime-representing function. \emph{Bull.
Amer. Math. Soc.} \textbf{53} (1947), 604. %
\url{https://doi.org/10.1090/S0002-9904-1947-08849-2}

\bibitem{Kuipers1950} L.~Kuipers. Prime-representing functions. \emph{Indag.
Math.} \textbf{12} (1950), 57--58. \url{https://zbmath.org/?q=an:0036.05102}

\bibitem{Ansari1951} A.~R.~Ansari. On prime representing function. \emph{%
Ganita} \textbf{2} (1951), 81--82. \url{https://zbmath.org/?q=an:0043.04302}

\bibitem{Wright1951} E.~M.~Wright. A prime-representing function. \emph{%
Amer. Math. Monthly} \textbf{58} (1951), 616--618. %
\url{https://doi.org/10.1080/00029890.1951.11999747}

\bibitem{Ore1952} O.~Ore. On the selection of subsequences. \emph{Proc.
Amer. Math. Soc.} \textbf{3} (1952), 706--712. %
\url{https://doi.org/10.1090/S0002-9939-1952-0050546-2}

\bibitem{Niven1951} I.~Niven. Functions which represent prime numbers. \emph{%
Proc. Amer. Math. Soc.} \textbf{2} (1951), 753--755. %
\url{https://doi.org/10.1090/S0002-9939-1951-0043807-4}

\bibitem{Wright1954} E.~M.~Wright. A class of representing functions. \emph{%
J. London Math. Soc.} \textbf{29} (1954), 63--71. %
\url{https://doi.org/10.1112/jlms/s1-29.1.63}

\bibitem{Ingham1937} A.~E.~Ingham. On the difference between consecutive
primes. \emph{Quart. J. Math. Oxford Ser.} (2) \textbf{8} (1937), 255--266. %
\url{https://doi.org/10.1093/qmath/os-8.1.255}

\bibitem{Dudley1969} U.~Dudley. History of a formula for primes. \emph{Amer.
Math. Monthly} \textbf{76} (1969), no.~1, 23--28. %
\url{https://doi.org/10.1080/00029890.1969.12000121}

\bibitem{HardyWright1960} G.~H.~Hardy and E.~M.~Wright. \emph{An
Introduction to the Theory of Numbers}, 4th ed. Oxford University Press,
1960. Theorem~418. %
\url{https://search.worldcat.org/search?q=An+Introduction+to+the+Theory+of+Numbers+Hardy+Wright}

\bibitem{Simpson2009} S.~G.~Simpson. \emph{Subsystems of Second Order
Arithmetic}, 2nd ed. Cambridge University Press, 2009. %
\url{https://doi.org/10.1017/CBO9780511581007}

\bibitem{FormalConjectures} Google DeepMind. \emph{formal-conjectures}. Lean
repository, 2026. \url{https://github.com/google-deepmind/formal-conjectures}

\bibitem{MathlibGraph} X.~Li and P.~Shafto. The Multinetwork of Mathlib.
NeurIPS 2026 Datasets and Benchmarks Track (under review). Dataset files
used for the dependency depth: \texttt{v2/declaration/metrics.csv} (column 
\texttt{dag\_layer} = topological depth of each declaration in the
declaration DAG) and \texttt{mathlib\_edges.csv} (raw dependency edges). 
{\scriptsize \ \url{https://huggingface.co/datasets/MathNetwork/MathlibGraph}%
.\newline
\url{https://huggingface.co/datasets/MathNetwork/MathlibGraph/resolve/main/v2/declaration/metrics.csv}%
\newline
\url{https://huggingface.co/datasets/MathNetwork/MathlibGraph/resolve/main/mathlib_edges.csv}%
.}

\bibitem{Bang1952} T.~Bang. A function representing all prime numbers. \emph{%
Nordisk Mat. Tidskr.} \textbf{34} (1952), 117--118. %
\url{https://zbmath.org/?q=an:0048.03108}

\bibitem{CaldwellCheng2005} C.~K.~Caldwell and Y.~Cheng. Determining Mills'
constant and a note on Honaker's problem. \emph{J. Integer Seq.} \textbf{8}
(2005), Article 05.4.1. %
\url{https://cs.uwaterloo.ca/journals/JIS/VOL8/Caldwell/caldwell78.pdf}

\bibitem{Haneke1962} W.~Haneke. Versch\"arfung der Absch\"atzung von $%
\zeta(1/2+it)$. \emph{Acta Arith.} \textbf{8} (1962/1963), 357--430. %
\url{https://doi.org/10.4064/aa-8-4-357-430}

\bibitem{Saito2025} K.~Saito. Mills' constant is irrational. \emph{%
Mathematika} \textbf{71} (2025), no.~3, e70027. %
\url{https://arxiv.org/abs/2404.19461}

\bibitem{Sierpinski1952} W.~Sierpi\'nski. Sur une formule donnant tous les
nombres premiers. \emph{C. R. Acad. Sci. Paris} \textbf{235} (1952),
1078--1079. \url{https://gallica.bnf.fr/ark:/12148/bpt6k31936/f1078}

\bibitem{Srinivasan1961} B.~R.~Srinivasan. Formulae for the $n$th prime. 
\emph{J. Indian Math. Soc.} (N.S.) \textbf{25} (1961), 33--39; \textbf{26}
(1962), 248. \url{https://zbmath.org/?q=an:0111.25002}

\bibitem{Srinivasan1962} B.~R.~Srinivasan. An arithmetical function and an
associated formula for the $n$th prime, I, II. \emph{Norske Vid. Selsk.
Forh. Trondheim} \textbf{35} (1962), 68--71, 72--75. %
\url{https://zbmath.org/?q=an:0117.03003}

\bibitem{Willans1964} C.~P.~Willans. On formulae for the $n$th prime number. 
\emph{Math. Gaz.} \textbf{48} (1964), 413--415. %
\url{https://doi.org/10.2307/3611815}
\end{thebibliography}

\end{document}
